% ============================================================
%  Stratified Monte Carlo Pricing and Variance Reduction
%  via K-Means Clustering and Machine Learning Control Variates
%
%  Benjamin Sulpice & Enzo Thuilliez
%  EDHEC Business School / Université Côte d'Azur
%  Supervised by Professor Arnaud DUFAYS
%  May 2026
%
%  Visual aesthetic: strictly monochrome/grayscale,
%  following the Pinelli & Gonçalves (2025) academic standard.
% ============================================================

\documentclass[12pt, a4paper]{report}

% ── Encoding & language ─────────────────────────────────────
\usepackage[T1]{fontenc}
\usepackage[utf8]{inputenc}
\usepackage[english]{babel}

% ── Typography: clean Computer Modern (lmodern) ─────────────
\usepackage{lmodern}

% ── Mathematics ─────────────────────────────────────────────
\usepackage{amsmath, amssymb, amsthm, mathtools}
\usepackage{bm}

% ── Page layout ─────────────────────────────────────────────
\setlength{\headheight}{15pt}
\usepackage[
  top=2.5cm, bottom=2.5cm,
  left=3cm,  right=2.5cm
]{geometry}
\usepackage{setspace}
\onehalfspacing

% ── Headers, titles, TOC ────────────────────────────────────
\usepackage{fancyhdr}
\usepackage{titlesec}
\usepackage{tocloft}
\usepackage{emptypage}

% ── Graphics ────────────────────────────────────────────────
\usepackage{graphicx}

% ── Tables ──────────────────────────────────────────────────
\usepackage{booktabs}
\usepackage{array}
\usepackage{caption}
\usepackage{subcaption}

% ── Floats ──────────────────────────────────────────────────
\usepackage{float}

% ── Code listings (for appendix) ────────────────────────────
\usepackage{listings}
\lstset{
  basicstyle=\ttfamily\small,
  keywordstyle=\bfseries,
  commentstyle=\itshape\color{gray},
  stringstyle=\ttfamily,
  showstringspaces=false,
  breaklines=true,
  frame=single,
  framerule=0.4pt,
  numbers=left,
  numberstyle=\tiny,
  numbersep=6pt,
  xleftmargin=1.5em,
  language=Python,
}

% ── Appendix ────────────────────────────────────────────────
\usepackage{appendix}

% ── Cross-references: hyperref with ALL colours set to black ─
\usepackage[
  hidelinks,
pdftitle = {Intro\textunderscore Research\textunderscore Thesis\textunderscore Sulpice\textunderscore Thuilliez.pdf},
pdfauthor  = {Benjamin Sulpice, Enzo Thuilliez},
  pdfsubject = {Mathematical Finance, Monte Carlo Methods, Machine Learning},
  pdfkeywords= {Monte Carlo, Variance Reduction, K-Means, Control Variates,
                Random Forest, Neural Network, Heston},
]{hyperref}
\usepackage{cleveref}

% ── Section title formatting — purely black, no colour ───────
\titleformat{\chapter}[display]
  {\normalfont\Large\bfseries}
  {\chaptertitlename\ \thechapter}{12pt}{\Huge}
\titlespacing*{\chapter}{0pt}{-20pt}{30pt}

\titleformat{\section}{\large\bfseries}{\thesection}{1em}{}
\titleformat{\subsection}{\normalsize\bfseries}{\thesubsection}{1em}{}

% ── Header / footer ─────────────────────────────────────────
\pagestyle{fancy}
\fancyhf{}
\fancyhead[L]{\small\textit{ML-Driven Variance Reduction in Monte Carlo Pricing}}
\fancyhead[R]{\small\textit{Sulpice \& Thuilliez}}
\fancyfoot[C]{\thepage}
\renewcommand{\headrulewidth}{0.4pt}

% ── Theorem environments ─────────────────────────────────────
\theoremstyle{definition}
\newtheorem{definition}{Definition}[chapter]
\newtheorem{proposition}{Proposition}[chapter]
\newtheorem{theorem}{Theorem}[chapter]
\newtheorem{remark}{Remark}[chapter]
\newtheorem{example}{Example}[chapter]

% ── Useful macros ────────────────────────────────────────────
\newcommand{\EQ}{\mathbb{E}^{\mathbb{Q}}}
\newcommand{\EP}{\mathbb{E}^{\mathbb{P}}}
\newcommand{\Q}{\mathbb{Q}}
\newcommand{\R}{\mathbb{R}}
\newcommand{\N}{\mathbb{N}}
\newcommand{\F}{\mathcal{F}}
\newcommand{\p}{\mathbb{P}}
\newcommand{\Var}{\operatorname{Var}}
\newcommand{\Cov}{\operatorname{Cov}}
\newcommand{\dd}{\,\mathrm{d}}

% ============================================================
\begin{document}

% ============================================================
% ── COVER PAGE ──────────────────────────────────────────────
% ============================================================

\begin{titlepage}

% ── LOGOS (Optically Balanced & Forced Single Line) ─────────
\begin{center}
    \noindent
    \begin{minipage}[b]{0.45\textwidth}
        \flushleft
        
\includegraphics[height=2.2cm, keepaspectratio]{logos/logo_EDHEC.png}
    \end{minipage}%
    \hfill
    \begin{minipage}[b]{0.45\textwidth}
        \flushright
        
\includegraphics[height=2.5cm, keepaspectratio]{logos/logo_UniCA.png}
    \end{minipage}
\end{center}

\vspace{1.5cm} % <── Increased space here to give the title breathing room

% ── TITLE BLOCK ─────────────────────────────────────────────
\begin{center}
    \rule{\textwidth}{0.6pt}
    \vspace{0.6cm}
    
    {\LARGE\bfseries Machine Learning-Driven Variance \\ Reduction in Monte Carlo Pricing\par}
    
    \vspace{0.5cm}
    {\large Stratified Clustering and Neural Network Control Variates\par}
    
    \vspace{0.6cm}
    \rule{\textwidth}{0.6pt}
\end{center}

\vspace{2.5cm}

% ── AUTHORS & SUPERVISOR ────────────────────────────────────
\begin{center}
    {\large\bfseries Benjamin Sulpice \quad \& \quad Enzo Thuilliez}
    
    \vspace{1.5cm}
    
    {\large Supervised by \textbf{Professor Arnaud DUFAYS}}
\end{center}

\vfill 

% ── DATE ────────────────────────────────────────────────────
\begin{center}
    {\large May 2026}
\end{center}

\end{titlepage}


% ── BLANK PAGE ──────────────────────────────────────────────
\cleardoublepage

\chapter*{Abstract}
\addcontentsline{toc}{chapter}{Abstract}
\thispagestyle{plain}

\noindent\textbf{Background and Context.}\enspace
Derivative pricing and market risk quantification occupy a central position
at the intersection of stochastic analysis, probability theory, and
computational mathematics.  Since Black, Scholes, and Merton formalised the
risk-neutral pricing framework in 1973, Monte Carlo simulation has emerged
as the method of choice for evaluating expectations under the risk-neutral
measure $\mathbb{Q}$, owing to its dimension-independence and natural
handling of path-dependent payoffs.

\medskip
\noindent\textbf{Core Problem Statement.}\enspace
The canonical Monte Carlo estimator converges at the slow rate
$\mathcal{O}(N^{-1/2})$, where $N$ denotes the number of simulated paths.
Halving the statistical error requires quadrupling the computational
budget, an unacceptable trade-off in production environments where prices
and risk sensitivities must be computed in near-real time.  Classical
variance reduction techniques (antithetic variates, control variates,
importance sampling) partially address this inefficiency but require
problem-specific design that limits their applicability across diverse
payoff structures.

\medskip
\noindent\textbf{Methodology.}\enspace
This memoir develops a unified machine learning framework for Monte Carlo
variance reduction.  Part~I establishes the mathematical foundations:
risk-neutral pricing under the Black--Scholes and Heston models and the
theoretical basis of classical variance reduction.  Part~II introduces
trajectory clustering (K-means and Gaussian Mixture Models) to partition
the simulation space into homogeneous regimes, then constructs data-driven
control variates using Random Forests and feedforward neural networks
trained to approximate the conditional payoff expectation
$\mathbb{E}^{\mathbb{Q}}[\Phi \mid \mathbf{x}]$.

\medskip
\noindent\textbf{Primary Findings.}\enspace
Clustering decomposes total payoff variance into between-cluster and
within-cluster components, attacking both through stratification and local
control variates respectively.  Random Forest control variates achieve
variance reduction factors of $10$--$50\times$ for standard European,
Asian, and barrier options while requiring no domain expertise in control
selection.  Neural network control variates extend these gains to
high-dimensional basket options and stochastic volatility settings,
achieving correlations $\rho > 0.90$ with the payoff on pilot samples of
$N_0 \geq 5 \times 10^4$ paths.

\medskip
\noindent\textbf{Key Implications.}\enspace
These results establish that learned variance reduction schemes match or
exceed classical hand-crafted methods at a fraction of the design effort,
with direct applicability to front-office pricing engines handling diverse
exotic payoffs.  The clustering-plus-local-learner pipeline generalises
seamlessly to new products without analytical rederivation, positioning
machine learning as a practical and principled component of production
Monte Carlo infrastructure.

\cleardoublepage

% ── TABLE OF CONTENTS ───────────────────────────────────────
\pagestyle{plain}
\setcounter{tocdepth}{2}

% TOC styling — monochrome
\renewcommand{\cftsecfont}{\normalsize}
\renewcommand{\cftchapleader}{\cftdotfill{\cftdotsep}}

\tableofcontents
\cleardoublepage

\pagestyle{fancy}

% ============================================================
% ── Introduction ────────────────────────────────────────────
% ============================================================
\chapter*{Introduction}
\addcontentsline{toc}{chapter}{Introduction}

The valuation of financial derivatives and the quantification of market
risk sit at the intersection of probability theory, stochastic analysis,
and numerical mathematics.  Over the past four decades, the rapid growth
of derivative markets has driven a remarkable convergence between
theoretical finance and applied mathematics, producing frameworks that
now underpin the daily operations of trading desks, risk management
units, and asset management firms worldwide.

At the core of modern derivative pricing lies a fundamental insight first
formalised by Black, Scholes, and Merton in 1973: under certain
no-arbitrage conditions, the price of a contingent claim can be expressed
as a discounted expectation under a \emph{risk-neutral} probability
measure.  For a broad class of payoff structures and underlying dynamics,
this expectation does not admit a closed-form solution, and one must
resort to numerical methods.  Among these, \textbf{Monte Carlo
simulation} occupies a privileged position: it is \emph{dimension-agnostic},
naturally handles path-dependent payoffs, and extends seamlessly to
multi-asset or multi-factor settings that defeat finite-difference or
lattice-based approaches.

The central drawback of Monte Carlo methods is their \emph{slow
convergence}: the statistical error in a Monte Carlo estimator decreases
only at rate $\mathcal{O}(N^{-1/2})$, where $N$ is the number of
simulated paths.  Halving the error requires \emph{quadrupling} the
computational budget, an unacceptable trade-off when prices and Greeks
must be computed in near-real time for hedging or client-quoting
purposes.  Variance reduction techniques seek to accelerate this
convergence without sacrificing the generality that makes Monte Carlo
attractive in the first place.

Classical variance reduction techniques such as antithetic variates,
control variates, and importance sampling require significant domain
expertise to be effective: the control must be chosen by hand, the change
of measure must be derived analytically, and optimal parameters must be
recomputed for each new product or model.  This lack of adaptability
limits their applicability in production environments handling hundreds
of distinct payoff structures.  Machine learning offers a fundamentally
different approach: rather than designing variance reduction schemes by
hand, one can learn them directly from simulated data, using flexible
function approximators that adapt automatically to the payoff and model
at hand, while preserving the martingale and no-arbitrage constraints
required for consistent pricing.

This memoir is organised around the following progression.  Part~I
establishes the mathematical foundations: the risk-neutral pricing
framework (Chapter~\ref{ch:pricing-mc}), and the theory of variance
reduction in its classical forms (Chapter~\ref{ch:variance-reduction}).
Part~II develops the machine learning reformulation, showing how optimal
control variates and importance sampling schemes can be \emph{learned}
from simulated data while preserving martingale and no-arbitrage
constraints.  Part~III provides a numerical study, implementing and
benchmarking the proposed methods on European, Asian, and barrier options
under both Black--Scholes and Heston dynamics.

Throughout the memoir, we privilege rigour without sacrificing economic
intuition.  Each mathematical development is immediately contextualised
by its practical implications for front-office pricing, hedging, and risk
management.

% ============================================================
% PART I — Mathematical Foundations
% ============================================================
\part{Mathematical Foundations}

% ── Chapter 1 ───────────────────────────────────────────────
\chapter{Derivative Pricing and Monte Carlo Simulation}
\label{ch:pricing-mc}

\section{The Risk-Neutral Pricing Framework}

\subsection{Probability spaces and filtrations}

We work on a filtered probability space
$(\Omega,\,\mathcal{F},\,(\mathcal{F}_t)_{t \in [0,T]},\,\mathbb{P})$,
where $\mathbb{P}$ is the \emph{physical} (or historical) probability
measure reflecting the actual dynamics of asset prices as observed in the
market.  The filtration $(\F_t)_{t \geq 0}$ encodes the information
available to market participants up to time $t$, and is assumed to
satisfy the usual conditions (right-continuity and completeness).

\begin{definition}[Contingent claim]
A \emph{contingent claim} with maturity $T > 0$ is an $\F_T$-measurable
random variable $\Phi \in L^2(\Omega, \F_T, \mathbb{P})$ representing
the payoff delivered at time $T$ contingent on the realised path of the
underlying.
\end{definition}

\begin{example}
For a European call option with strike $K$ on an underlying $S_T$, the
payoff is $\Phi = (S_T - K)^+$.  For an Asian call,
$\Phi = \bigl(\bar{S}_T - K\bigr)^+$ where
$\bar{S}_T = \frac{1}{T}\int_0^T S_t \dd t$ is the running average.
\end{example}

\subsection{No-arbitrage and the risk-neutral measure}

The \emph{First Fundamental Theorem of Asset Pricing} (Harrison \&
Kreps, 1979; Delbaen \& Schachermayer, 1994) establishes the equivalence
between the absence of arbitrage and the existence of an \emph{equivalent
martingale measure} (EMM).

\begin{theorem}[First Fundamental Theorem]
Let $r \geq 0$ be the continuously compounded risk-free rate and let
$\tilde{S}_t = e^{-rt}S_t$ denote the discounted price process.  The
market is free of arbitrage (in the sense of NFLVR) if and only if there
exists a probability measure $\Q \sim \mathbb{P}$ under which
$\tilde{S}$ is a local martingale.  Such a measure $\Q$ is called a
\emph{risk-neutral} (or \emph{pricing}) measure.
\end{theorem}

When the EMM is unique, the market is said to be \emph{complete}: every
contingent claim is replicable by a self-financing trading strategy.  The
Black--Scholes model, driven by a single Brownian motion, is a canonical
example of a complete market.  Models with stochastic volatility (e.g.\
Heston) or jumps introduce incompleteness and give rise to a family of
pricing measures.

\subsection{The fundamental pricing equation}

Under the risk-neutral measure $\Q$, the fair value at time $0$ of a
contingent claim $\Phi$ is given by the discounted risk-neutral
expectation:

\begin{equation}
  \boxed{V_0 = e^{-rT}\,\EQ\!\left[\Phi\right]}
  \label{eq:pricing}
\end{equation}

\noindent where the expectation is taken with respect to $\Q$.  More
generally, the price process $(V_t)_{t \in [0,T]}$ is a
$\Q$-martingale after discounting, satisfying:
\[
  V_t = e^{-r(T-t)}\,\EQ\!\left[\Phi \,\middle|\, \F_t\right],
  \qquad t \in [0,T].
\]

Equation~\eqref{eq:pricing} is the starting point for all numerical
pricing methods.  The challenge lies in computing this expectation
efficiently when no analytical formula is available.

\section{The Black--Scholes Model}

\subsection{Geometric Brownian motion}

The \emph{Black--Scholes--Merton} (BSM) model (Black \& Scholes, 1973;
Merton, 1973) postulates that the underlying asset price follows a
\emph{geometric Brownian motion} (GBM) under the physical measure:
\[
  \dd S_t = \mu\, S_t\, \dd t + \sigma\, S_t\, \dd W_t^{\mathbb{P}},
\]
where $\mu \in \R$ is the drift, $\sigma > 0$ is the (constant)
volatility, and $W^{\mathbb{P}}$ is a standard Brownian motion under
$\mathbb{P}$.

By Girsanov's theorem, there exists a unique EMM $\Q$ under which
$W_t^{\Q} = W_t^{\mathbb{P}} + \frac{\mu - r}{\sigma}t$ is a standard
Brownian motion (the Sharpe ratio $\frac{\mu-r}{\sigma}$ is the
\emph{market price of risk}).  Under $\Q$:
\[
  \dd S_t = r\, S_t\, \dd t + \sigma\, S_t\, \dd W_t^{\Q}.
\]

Applying Itô's lemma to $\log S_t$ yields the closed-form solution:
\begin{equation}
  S_T = S_0 \exp\!\left[\left(r - \tfrac{1}{2}\sigma^2\right)T
      + \sigma\, W_T^{\Q}\right]
      = S_0 \exp\!\left[\left(r - \tfrac{1}{2}\sigma^2\right)T
      + \sigma\sqrt{T}\,Z\right],
  \label{eq:gbm-solution}
\end{equation}
where $Z \sim \mathcal{N}(0,1)$ under $\Q$.

\subsection{The Black--Scholes formula}

For a European call option $\Phi = (S_T - K)^+$, inserting
\eqref{eq:gbm-solution} into \eqref{eq:pricing} and evaluating the
Gaussian integral yields the celebrated closed-form:

\begin{equation}
  V_0^{\text{call}} = S_0\,\mathcal{N}(d_1) - K e^{-rT}\,\mathcal{N}(d_2),
  \label{eq:BS-call}
\end{equation}
where
\[
  d_1 = \frac{\ln(S_0/K) + (r + \sigma^2/2)T}{\sigma\sqrt{T}}, \qquad
  d_2 = d_1 - \sigma\sqrt{T},
\]
and $\mathcal{N}$ denotes the standard normal CDF.  Put-call parity
immediately gives the put price.

Despite its analytical elegance, the BSM model carries well-documented
empirical shortcomings: the volatility surface observed in liquid options
markets exhibits a pronounced \emph{skew} and \emph{smile} structure
inconsistent with the constant-$\sigma$ assumption.  Path-dependent or
exotic payoffs also render \eqref{eq:BS-call} inapplicable, motivating
the need for simulation-based pricing.

\section{The Heston Stochastic Volatility Model}

\subsection{Model specification}

To address the empirical failures of BSM, Heston (1993) proposed a model
in which volatility is itself a stochastic process.  Under $\Q$, the
joint dynamics are:
\begin{align}
  \dd S_t &= r\,S_t\,\dd t + \sqrt{V_t}\,S_t\,\dd W_t^S, \label{eq:heston-S}\\
  \dd V_t &= \kappa(\theta - V_t)\,\dd t + \xi\sqrt{V_t}\,\dd W_t^V, \label{eq:heston-V}\\
  \dd\langle W^S, W^V\rangle_t &= \rho\,\dd t, \label{eq:heston-corr}
\end{align}
where $V_t = \sigma_t^2$ is the instantaneous variance, $\kappa > 0$ is
the mean-reversion speed, $\theta > 0$ is the long-run variance,
$\xi > 0$ is the volatility of volatility, and $\rho \in (-1, 1)$ is
the correlation between the two Brownian motions.

\begin{definition}[Feller condition]
The \emph{Feller condition} $2\kappa\theta > \xi^2$ ensures that the
variance process $V_t$ remains strictly positive almost surely,
guaranteeing the well-posedness of the model.
\end{definition}

The Heston model can reproduce the implied volatility smile through the
interplay of $\rho$ (controlling skew) and $\xi$ (controlling the
curvature of the smile).  Semi-analytical prices for vanilla options are
available via Fourier inversion of the characteristic function of
$\log S_T$, but exotic payoffs again require Monte Carlo simulation.

\subsection{Euler--Maruyama discretisation}

For simulation purposes, the continuous-time
system~\eqref{eq:heston-S}--\eqref{eq:heston-V} is discretised on a
uniform time grid $0 = t_0 < t_1 < \cdots < t_M = T$ with step
$\Delta t = T/M$.  The \emph{Euler--Maruyama} scheme reads:
\begin{align*}
  V_{t_{k+1}} &= V_{t_k} + \kappa(\theta - V_{t_k})\Delta t
                 + \xi\sqrt{V_{t_k}^+}\,\Delta W^V_k,\\
  S_{t_{k+1}} &= S_{t_k}\exp\!\left[\left(r - \tfrac{1}{2}V_{t_k}^+\right)\Delta t
                 + \sqrt{V_{t_k}^+\,\Delta t}\,Z_k^S\right],
\end{align*}
where $V^+ = \max(V,0)$ is the \emph{full truncation} scheme to prevent
negative variance, $\Delta W^V_k \sim \mathcal{N}(0,\Delta t)$, and
$(Z_k^S, \Delta W^V_k)$ are drawn from a bivariate normal with
correlation $\rho$.

More sophisticated discretisation schemes --- Milstein, QE (Andersen,
2008), BK --- offer superior bias properties at the cost of additional
complexity; these are discussed in Part~III.

\section{Monte Carlo Simulation for Derivative Pricing}

\subsection{The Monte Carlo estimator}

The Monte Carlo approach to evaluating \eqref{eq:pricing} is
conceptually straightforward.  One simulates $N$ independent paths of
the underlying under $\Q$, computes the payoff on each path, and takes
the discounted sample mean:

\begin{equation}
  \hat{V}_0^{MC} = e^{-rT}\,\frac{1}{N}\sum_{i=1}^{N} \Phi(\omega_i),
  \label{eq:mc-estimator}
\end{equation}
where $\omega_i$ denotes the $i$-th simulated path.  By the Strong Law
of Large Numbers, $\hat{V}_0^{MC} \to V_0$ almost surely as
$N \to \infty$, provided $\EQ[|\Phi|] < \infty$.

\begin{figure}[htbp]
  \centering
  \includegraphics[width=0.65\textwidth]{figures/figure_1_trajectoires.png}
  \caption{Two hundred simulated trajectories of the underlying asset
           price $S_t$ under the risk-neutral measure $\mathbb{Q}$,
           following geometric Brownian motion with parameters
           $S_0 = K = 100$, $r = 5\%$, $\sigma = 20\%$, $T = 1$ year
           and $M = 252$ daily time steps.
           \textit{Left:} European call --- paths finishing above the
           strike $K$ (in-the-money, dark blue) contribute a strictly
           positive payoff $(S_T - K)^+$, while paths finishing below
           (light blue) yield a zero payoff.
           \textit{Right:} European put --- in-the-money paths (dark red)
           finish below $K$.
           The dashed line marks the strike level; the Black--Scholes
           analytical price $V_0$ is reported for each contract.}
  \label{fig:mc-trajectories}
\end{figure}

\subsection{Convergence and the $\mathcal{O}(N^{-1/2})$ rate}

The quality of the estimator~\eqref{eq:mc-estimator} is measured by its
\emph{root mean squared error} (RMSE), which by the Central Limit
Theorem satisfies:

\begin{equation}
  \mathrm{RMSE}(\hat{V}_0^{MC}) = \frac{e^{-rT}\,\sigma_\Phi}{\sqrt{N}}
  = \mathcal{O}\!\left(N^{-1/2}\right),
  \label{eq:mc-convergence}
\end{equation}
where $\sigma_\Phi^2 = \Var_{\Q}[\Phi]$ is the variance of the
(undiscounted) payoff, as illustrated in Figure~\ref{fig:mc-rmse}.
Crucially, this rate is \emph{dimension-independent}: unlike
finite-difference methods, which suffer from the \emph{curse of
dimensionality} (error scaling as $\mathcal{O}(h^p)$ where $h$ is the
grid spacing in each of $d$ dimensions), the Monte Carlo rate
$\mathcal{O}(N^{-1/2})$ is unaffected by the number of risk factors.

The $\mathcal{O}(N^{-1/2})$ rate implies that achieving double
precision (halving the error) requires \emph{four times} as many
simulations.  For reference, a typical front-office pricing engine
running $N = 10^6$ paths achieves a relative error of roughly $0.1\%$
--- acceptable for vanilla options but insufficient for Delta/Vega
hedging where Greeks are obtained by finite differences and the noise is
amplified.

\begin{figure}[H]
  \centering
  \includegraphics[width=0.85\textwidth]{figures/figure_2_RMSE.png}
  \caption{Empirical RMSE of the Monte Carlo estimator $\hat{V}_0^{MC}$
           as a function of the number of simulated paths $N$, for a
           European call (solid blue) and a European put (dashed red),
           under Black--Scholes dynamics with parameters $S_0 = K = 100$,
           $r = 5\%$, $\sigma = 20\%$, $T = 1$ year.  The RMSE is
           estimated over $200$ independent replications for each value
           of $N$.  The dotted black line represents the theoretical
           $\mathcal{O}(N^{-1/2})$ reference slope.  Both axes are on
           a logarithmic scale, confirming the expected decay rate:
           halving the RMSE requires quadrupling the computational
           budget.}
  \label{fig:mc-rmse}
\end{figure}

\subsection{The variance problem and its consequences}

The payoff variance $\sigma_\Phi^2$ is the sole driver of the Monte
Carlo error in~\eqref{eq:mc-convergence}.  For several practically
important payoff types, this variance is large:

\begin{itemize}
  \item \textbf{Deep out-of-the-money options}: the payoff is zero with
        high probability and very large on rare events, leading to
        heavy-tailed distributions and extreme variance.
  \item \textbf{Digital (binary) options}: the payoff jumps
        discontinuously, producing high sensitivity to the exact
        crossing of a threshold.
  \item \textbf{Barrier options}: the payoff depends on whether a path
        crosses a barrier level, introducing additional path-dependency
        that inflates variance when the barrier is close to the spot.
  \item \textbf{Greeks via finite differences}: even when the option
        price is estimated with acceptable accuracy, the
        finite-difference Greek estimator
        $\hat{\Delta} \approx (\hat{V}_0(S_0 + \varepsilon) -
        \hat{V}_0(S_0))/\varepsilon$ can have catastrophically high
        variance as $\varepsilon \to 0$.
\end{itemize}

These limitations motivate the study of \emph{variance reduction
techniques}, which aim to decrease $\sigma_\Phi^2$ --- or equivalently
increase the effective sample size $N$ for a fixed computational budget
--- without introducing bias into the estimator.

\begin{remark}
A complementary approach to variance reduction is \emph{Quasi-Monte
Carlo} (QMC), which replaces pseudo-random draws with deterministic
low-discrepancy sequences (e.g.\ Sobol', Halton) to achieve convergence
rates approaching $\mathcal{O}(N^{-1})$ in practice.  While powerful,
QMC is less straightforward to combine with adaptive machine learning
schemes; we briefly revisit it in Chapter~\ref{ch:variance-reduction}.
\end{remark}

\subsection{Confidence intervals and practical diagnostics}

For a given simulation with $N$ paths, a $(1-\alpha)$ confidence
interval for $V_0$ is:
\[
  \hat{V}_0^{MC} \pm z_{\alpha/2}\,\frac{e^{-rT}\hat{\sigma}_\Phi}{\sqrt{N}},
\]
where $\hat{\sigma}_\Phi$ is the sample standard deviation of the
payoffs and $z_{\alpha/2}$ is the appropriate normal quantile
($z_{0.025} = 1.96$ for $95\%$ coverage).

In practice, one monitors the \emph{coefficient of variation}
$\mathrm{CV} = \hat{\sigma}_\Phi / |\hat{V}_0^{MC}|$ as a relative
measure of simulation noise.  A CV above $10\%$ is typically
unacceptable in a production environment.

\section{Summary}

This chapter has established the conceptual and mathematical foundations
for the rest of the memoir.  We have:

\begin{enumerate}
  \item Introduced the risk-neutral pricing framework and derived the
        fundamental pricing equation $V_0 = e^{-rT}\EQ[\Phi]$.
  \item Presented the Black--Scholes model, its closed-form solution for
        vanilla options, and the Euler--Maruyama simulation scheme.
  \item Extended the framework to the Heston stochastic volatility model,
        motivating the need for path simulation.
  \item Formalised the Monte Carlo estimator and its
        $\mathcal{O}(N^{-1/2})$ convergence rate, and identified the
        high-variance regimes that limit its practical applicability.
\end{enumerate}

Chapter~\ref{ch:variance-reduction} now turns to classical variance
reduction techniques --- antithetic variables, control variates, and
importance sampling --- which form the conceptual precursors to the
machine-learning methods developed in Part~II.

% ============================================================
%  CHAPTER 2 -- Classical Variance Reduction Techniques
% ============================================================
\chapter{Classical Variance Reduction Techniques}
\label{ch:variance-reduction}

As established in Chapter~\ref{ch:pricing-mc}, the Monte Carlo
estimator converges at the slow rate $\mathcal{O}(N^{-1/2})$, with
the error driven entirely by the payoff variance
$\sigma_\Phi^2 = \Var_{\Q}[\Phi]$.  The natural question is therefore:
can one construct an alternative unbiased estimator with strictly smaller
variance, without increasing the computational budget?  This chapter
presents four classical methods that answer affirmatively ---
antithetic variates, control variates, importance sampling, and
stratified sampling --- each exploiting a different structural property
of the pricing problem.  Together, they form the theoretical baseline
against which the machine learning methods of Part~II are evaluated.

\section{Antithetic Variates}
\label{sec:antithetic}

The antithetic variates method is the simplest variance reduction
technique.  It rests on a single observation: for any integrable function
$\Phi$ and a symmetric random variable $Z \sim \mathcal{N}(0,1)$, the
pair $(\Phi(Z),\, \Phi(-Z))$ has the same marginal distribution as
$(\Phi(Z),\, \Phi(Z))$, but may be negatively correlated.  Averaging the
two gives a lower-variance estimator at no additional simulation cost.

\begin{definition}[Antithetic estimator]
Let $Z_1, \ldots, Z_N \overset{\text{i.i.d.}}{\sim} \mathcal{N}(0,1)$.
The antithetic estimator is:
\begin{equation}
    \hat{V}_0^{AV} = \frac{e^{-rT}}{N} \sum_{i=1}^{N}
    \frac{\Phi(Z_i) + \Phi(-Z_i)}{2}.
    \label{eq:av-estimator}
\end{equation}
\end{definition}

The estimator is unbiased by symmetry of the normal distribution.  Its
variance satisfies:
\begin{equation}
    \Var\!\left[\hat{V}_0^{AV}\right]
    = \frac{e^{-2rT}}{2N}
    \Bigl(\Var[\Phi(Z)] + \Cov[\Phi(Z),\,\Phi(-Z)]\Bigr).
    \label{eq:av-variance}
\end{equation}

Variance reduction occurs whenever $\Cov[\Phi(Z),\Phi(-Z)] < 0$, which
is guaranteed when $\Phi$ is a \emph{monotone} function of $Z$.  For a
European call under Black--Scholes, the payoff
$(S_0 e^{(r - \sigma^2/2)T + \sigma\sqrt{T}Z} - K)^+$ is strictly
increasing in $Z$, so antithetic pairs are always negatively correlated
in this setting.  Empirically, the method yields variance reductions of
$5$--$10\times$ for at-the-money vanilla options (Glasserman, 2004).

The method is less effective for non-monotone payoffs such as straddles
or certain path-dependent structures, where
$\Cov[\Phi(Z), \Phi(-Z)]$ may be close to zero or even positive.

\section{Control Variates}
\label{sec:control-variates}

Control variates exploit a correlated quantity whose expectation is known
analytically.  Let $Y$ be a random variable --- the \emph{control} ---
with known mean $\mu_Y = \mathbb{E}^{\Q}[Y]$ and positively correlated
with $\Phi$.  For any scalar $c \in \R$, the estimator:
\begin{equation}
    \hat{V}_0^{CV}(c) = \frac{e^{-rT}}{N} \sum_{i=1}^{N}
    \bigl[\Phi(\omega_i) - c\,(Y_i - \mu_Y)\bigr]
    \label{eq:cv-estimator}
\end{equation}
is unbiased for all $c$, since $\mathbb{E}[Y_i - \mu_Y] = 0$.

\begin{proposition}[Optimal coefficient]
The variance of $\hat{V}_0^{CV}(c)$ is minimised by:
\begin{equation}
    c^* = \frac{\Cov[\Phi,\, Y]}{\Var[Y]},
    \label{eq:cv-optimal}
\end{equation}
yielding the reduced variance:
\begin{equation}
    \Var\!\left[\hat{V}_0^{CV}(c^*)\right]
    = \Var[\Phi]\,\bigl(1 - \rho_{\Phi,Y}^2\bigr),
    \label{eq:cv-variance}
\end{equation}
where $\rho_{\Phi,Y} = \mathrm{Corr}[\Phi,\, Y]$.
\end{proposition}

The factor $(1 - \rho_{\Phi,Y}^2)$ quantifies the gain: a correlation
of $|\rho| = 0.95$ already reduces variance by a factor of $10$.  The
challenge lies in identifying a useful control.  In the Black--Scholes
setting, the geometric Asian call --- whose price is available in closed
form (Kemna \& Vorst, 1990) --- serves as a natural control for the
arithmetic Asian call, typically achieving $|\rho| > 0.99$.

In practice, $c^*$ is unknown and estimated from the same simulation
batch via ordinary least squares regression of $\Phi$ on $Y$,
introducing only a negligible finite-sample bias for moderate $N$.

\section{Importance Sampling}
\label{sec:importance-sampling}

Importance sampling addresses a different source of inefficiency: for
out-of-the-money options, most simulated paths yield a zero payoff,
contributing nothing to the estimator while consuming computational
resources.  The idea is to simulate under an alternative measure
$\tilde{\Q}$ that concentrates paths in the payoff-relevant region, and
to correct for the change of measure via a likelihood ratio.

Let $\tilde{\Q}$ be a probability measure equivalent to $\Q$, with
Radon--Nikodym derivative $L = \dd\Q / \dd\tilde{\Q}$.  Then:
\begin{equation}
    V_0 = e^{-rT}\,\EQ[\Phi]
        = e^{-rT}\,\mathbb{E}^{\tilde{\Q}}[\Phi \cdot L],
    \label{eq:is-identity}
\end{equation}
which motivates the importance sampling estimator:
\begin{equation}
    \hat{V}_0^{IS} = \frac{e^{-rT}}{N} \sum_{i=1}^{N}
    \Phi(\tilde{\omega}_i)\,L(\tilde{\omega}_i),
    \label{eq:is-estimator}
\end{equation}
where $\tilde{\omega}_1, \ldots, \tilde{\omega}_N \sim \tilde{\Q}$.
The estimator is unbiased by~\eqref{eq:is-identity}.

\begin{proposition}[Zero-variance measure]
The variance of $\hat{V}_0^{IS}$ is zero if and only if:
\[
    \frac{\dd\tilde{\Q}^*}{\dd\Q} \propto \Phi.
\]
\end{proposition}

This result is theoretically exact but practically circular, since
constructing $\tilde{\Q}^*$ requires knowing $V_0$.  In the
Black--Scholes setting, a tractable approximation is \emph{exponential
tilting}: one shifts the drift of $W^{\Q}$ by a constant $\theta$, so
that the mean of $\log S_T$ under $\tilde{\Q}_\theta$ is shifted toward
$\log K$.  The likelihood ratio takes the simple form
$L_\theta = \exp(-\theta W_T - \frac{1}{2}\theta^2 T)$, and the
optimal $\theta$ is found by minimising the variance of
$\Phi \cdot L_\theta$ numerically.  For deep out-of-the-money options,
this approach can reduce variance by several orders of magnitude.

In high-dimensional settings, however, designing effective changes of
measure becomes significantly harder: poorly chosen $\tilde{\Q}$ leads
to \emph{weight degeneracy}, where a few paths carry extreme likelihood
ratios, inflating the variance of the weighted estimator.

\section{Stratified Sampling}
\label{sec:stratification}

Stratified sampling ensures that the simulation covers the probability
space more uniformly than pure random sampling.  One partitions the
driving uniform variable $U \in [0,1]$ into $K$ equal-probability
strata $A_k = \bigl[\frac{k-1}{K}, \frac{k}{K}\bigr)$, draws one
sample per stratum, and uses the inverse transform to obtain the
corresponding path.  The resulting estimator is:
\begin{equation}
    \hat{V}_0^{ST} = \frac{e^{-rT}}{N} \sum_{k=1}^{K}
    \sum_{i=1}^{n_k} \Phi\!\left(F^{-1}\!\left(U_i^{(k)}\right)\right),
    \label{eq:strat-estimator}
\end{equation}
where $U_i^{(k)} \sim \mathrm{Uniform}(A_k)$ and $F^{-1}$ is the
quantile function of the relevant marginal distribution.

The variance reduction relative to crude Monte Carlo is:
\begin{equation}
    \Var[\hat{V}_0^{MC}] - \Var[\hat{V}_0^{ST}]
    = \frac{e^{-2rT}}{N} \sum_{k=1}^{K} p_k\,(\mu_k - V_0 e^{rT})^2
    \geq 0,
    \label{eq:strat-gain}
\end{equation}
where $\mu_k = \EQ[\Phi \,|\, A_k]$ is the within-stratum conditional
mean.  The gain is large when the stratum means differ --- i.e.\ when
the chosen stratification variable is informative about the payoff.

For a fixed total budget $N$, variance is further minimised by the
\emph{Neyman allocation} $n_k^* \propto p_k \sigma_k$, which assigns
more paths to strata with higher within-stratum standard deviation
$\sigma_k$.

\section{Summary and Motivation for Machine Learning}
\label{sec:vr-summary}

% Table uses booktabs (no vertical lines) — upgrade from original
\begin{table}[H]
\centering
\caption{Comparative overview of classical variance reduction
         techniques.  Implementation difficulty and high-dimension
         performance are qualitative assessments; typical variance
         gains are reported from Glasserman (2004).}
\label{tab:vr-comparison}
\renewcommand{\arraystretch}{1.4}
\begin{tabular}{lcccc}
\toprule
\textbf{Method} & \textbf{Unbiased} & \textbf{Implementation}
                & \textbf{Typical gain} & \textbf{High dimension} \\
\midrule
Antithetic variates  & Yes & Trivial  & $2$--$10\times$    & Poor \\
Control variates     & Yes & Easy     & $10$--$100\times$  & Moderate \\
Importance sampling  & Yes & Moderate & up to $1000\times$ & Poor \\
Stratified sampling  & Yes & Moderate & $2$--$20\times$    & Moderate \\
\bottomrule
\end{tabular}
\end{table}

All four methods share a common limitation: they require
\emph{problem-specific design}.  Antithetic variates require a monotone
payoff; control variates require an analytically tractable correlated
quantity; importance sampling requires a tractable approximate change of
measure; stratification requires a well-chosen partition variable.  Each
optimal parameter --- $c^*$, $\theta^*$, $n_k^*$ --- must be re-derived
whenever the model or payoff changes, which is impractical in a
front-office environment handling hundreds of distinct products.

This motivates the central contribution of Part~II: rather than
designing variance reduction schemes by hand, one can \emph{learn} them
directly from simulated data.  Specifically, neural networks can
parametrise both the control variate function and the change-of-measure
density, trained to minimise the empirical variance of the Monte Carlo
estimator subject to the unbiasedness constraint.  The classical methods
of this chapter then become special cases of a general learned framework,
and the theoretical upper bounds on variance reduction they establish
serve as benchmarks for the machine learning approach.

% ============================================================
%  PART II — Machine Learning for Variance Reduction
% ============================================================
\part{Machine Learning for Variance Reduction}

% ── Chapter 3 ───────────────────────────────────────────────
\chapter{Clustering of Monte Carlo Trajectories}
\label{ch:clustering}

The classical variance reduction techniques of
Chapter~\ref{ch:variance-reduction} apply a single global estimator
uniformly across all simulated paths.  This is suboptimal when the
simulation space is heterogeneous: a path that ends deep in-the-money
and a path that ends deep out-of-the-money carry very different
information about the option price, yet they are treated identically by
a global control variate or a global importance sampling scheme.

The idea of this chapter is simple: before applying any variance
reduction, \emph{group} the simulated paths into clusters of similar
behaviour, and then apply variance reduction \emph{locally} within each
cluster.  Paths within the same cluster are more homogeneous than the
full population, so local estimators are more accurate.  This is the
core intuition behind trajectory clustering for Monte Carlo simulation.

The chapter is organised as follows.  Section~\ref{sec:features}
explains how to summarise each path by a small set of numbers (its
\emph{features}), and why normalising those features is essential.
Section~\ref{sec:kmeans} introduces the K-means clustering algorithm.
Section~\ref{sec:gmm} presents Gaussian Mixture Models as a more
flexible alternative.  Sections~\ref{sec:adaptive-allocation}
and~\ref{sec:local-cv} then show precisely how clustering reduces
variance, through adaptive allocation and local control variates
respectively.

\section{From Paths to Feature Vectors}
\label{sec:features}

\subsection{Why not cluster paths directly?}

A Monte Carlo path
$\omega_i = (S_{t_0}^{(i)}, \ldots, S_{t_M}^{(i)})$ with $M = 252$
daily steps lives in $\R^{252}$.  Clustering directly in this space is
problematic for two reasons.  First, computing distances between
$N = 10^5$ vectors in $\R^{252}$ is computationally expensive.  Second,
and more fundamentally, distance measures become unreliable in high
dimensions: in $\R^{252}$, most pairs of points are approximately
equidistant from each other, making it impossible to define meaningful
clusters.  This phenomenon is known as the \emph{curse of
dimensionality}.

The solution is to map each path to a small number $d$ of summary
statistics --- its \emph{feature vector}
$\mathbf{x}_i \in \R^d$ with $d \ll M$ --- that capture the financially
relevant behaviour of the path.  Clustering is then performed in this
low-dimensional feature space.

\subsection{Financially motivated features}

The choice of features is guided by the type of option being priced.
The following five statistics cover the most common situations.

\textbf{Terminal value} $x_i^{(1)} = S_T^{(i)}$.  For European options,
the payoff depends only on $S_T$, so this single number is the most
informative feature.

\textbf{Realised volatility.}  The annualised realised variance measures
how much the path fluctuated along the way:
\begin{equation}
    x_i^{(2)} = \hat{\sigma}_i^2
    = \frac{1}{(M-1)\Delta t} \sum_{k=1}^{M}
    \left(\log\frac{S_{t_k}^{(i)}}{S_{t_{k-1}}^{(i)}}\right)^2.
    \label{eq:realised-vol}
\end{equation}

\textbf{Maximum drawdown.}  The largest peak-to-trough decline:
\begin{equation}
    x_i^{(3)} = \max_{0 \leq j \leq k \leq M}
    \frac{S_{t_j}^{(i)} - S_{t_k}^{(i)}}{S_{t_j}^{(i)}},
    \label{eq:max-drawdown}
\end{equation}
is particularly relevant for barrier and lookback options.

\textbf{Barrier hitting time.}  For a barrier option with barrier level
$B$:
\begin{equation}
    x_i^{(4)} = \tau_B^{(i)}
    = \min\bigl\{t_k : S_{t_k}^{(i)} \leq B\bigr\},
    \label{eq:hitting-time}
\end{equation}
(set to $T$ if the barrier is never hit).

\textbf{Running average.}  For Asian options:
\[
    x_i^{(5)} = \bar{S}^{(i)} = \frac{1}{M}\sum_{k=1}^{M} S_{t_k}^{(i)}.
\]

In practice, one selects $d = 2$ to $5$ features depending on the
option type, and proceeds to the normalisation step below.

\subsection{Feature normalisation}

Before clustering, it is essential to \emph{normalise} the features.
The standard remedy is \emph{standardisation}: for each feature
$\ell \in \{1, \ldots, d\}$, replace each raw value by:
\begin{equation}
    \tilde{x}_i^{(\ell)}
    = \frac{x_i^{(\ell)} - \bar{x}^{(\ell)}}{\hat{s}^{(\ell)}},
    \label{eq:standardisation}
\end{equation}
where $\bar{x}^{(\ell)}$ and $\hat{s}^{(\ell)}$ are the sample mean
and standard deviation of feature $\ell$ across all $N$ paths.

\section{K-Means Clustering}
\label{sec:kmeans}

\subsection{The core idea}

The goal of clustering is to partition the $N$ feature vectors
$\mathbf{x}_1, \ldots, \mathbf{x}_N \in \R^d$ into $K$ groups such
that points within the same group are as similar as possible.  K-means
formalises this by minimising the total squared distance from each point
to the centre of its group.

To measure similarity, K-means uses the \emph{Euclidean distance}:
\begin{equation}
    d(\mathbf{x}_i, \mathbf{x}_j)
    = \|\mathbf{x}_i - \mathbf{x}_j\|_2
    = \sqrt{\sum_{\ell=1}^{d}
      \bigl(x_i^{(\ell)} - x_j^{(\ell)}\bigr)^2}.
    \label{eq:euclidean}
\end{equation}

\subsection{Objective and algorithm}

Given $K$ clusters $C_1, \ldots, C_K$ and their centroids
$\boldsymbol{\mu}_1, \ldots, \boldsymbol{\mu}_K \in \R^d$, the
\emph{within-cluster inertia} measures how compact the clusters are:
\begin{equation}
    \mathcal{J}(\mathcal{C}) = \sum_{k=1}^{K} \sum_{i \in C_k}
    \|\mathbf{x}_i - \boldsymbol{\mu}_k\|^2.
    \label{eq:kmeans-objective}
\end{equation}

\begin{definition}[K-means algorithm]
Starting from $K$ randomly initialised centroids:
\begin{enumerate}
    \item \textbf{Assignment}: assign each point to its nearest centroid,
    \[
        k^*(i) = \arg\min_{k \in \{1,\ldots,K\}}
        \|\mathbf{x}_i - \boldsymbol{\mu}_k\|^2;
    \]
    \item \textbf{Update}: recompute each centroid as the mean of its
    assigned points,
    \[
        \boldsymbol{\mu}_k
        = \frac{1}{|C_k|}\sum_{i \in C_k} \mathbf{x}_i.
    \]
\end{enumerate}
These two steps repeat until no point changes cluster between two
consecutive iterations.
\end{definition}

\subsection{Limitations of K-means}
\label{sec:kmeans-limitations}

Three structural limitations motivate the introduction of Gaussian
Mixture Models.  First, K-means implicitly assumes that clusters are
spherical and roughly equal in size; financial path features are often
elongated or asymmetric.  Second, each path belongs to exactly one
cluster with no measure of confidence, inflating within-cluster variance
near boundaries.  Third, the centroid is the arithmetic mean and is
sensitive to extreme values.  The K-means++ initialisation (Arthur \&
Vassilvitskii, 2007) partially addresses the third issue.

\subsection{Choosing the number of clusters $K$}

Two standard methods identify a good value of $K$.

\textbf{Elbow method.}  Plot the inertia $\mathcal{J}(K)$ as a
function of $K$ and identify the value beyond which the gain becomes
small.

\textbf{Silhouette coefficient.}  For each point $\mathbf{x}_i$ in
cluster $C_k$, define $a_i$ as the mean within-cluster distance and
$b_i$ as the mean distance to the nearest other cluster.  The
silhouette score is:
\begin{equation}
    s_i = \frac{b_i - a_i}{\max(a_i,\, b_i)} \in [-1, 1].
    \label{eq:silhouette}
\end{equation}
The optimal $K$ maximises the average silhouette
$\bar{s}(K) = \frac{1}{N}\sum_i s_i$, with $\bar{s} > 0.5$ indicating
well-separated clusters.

\section{Gaussian Mixture Models}
\label{sec:gmm}

\subsection{Model definition}

\begin{definition}[Gaussian Mixture Model]
A GMM with $K$ components models the density of $\mathbf{x} \in \R^d$
as:
\begin{equation}
    p(\mathbf{x}) = \sum_{k=1}^{K} \pi_k\,
    \mathcal{N}(\mathbf{x};\, \boldsymbol{\mu}_k,\,
    \boldsymbol{\Sigma}_k),
    \label{eq:gmm}
\end{equation}
where $\pi_k > 0$ with $\sum_k \pi_k = 1$ are the \emph{mixing
weights}, $\boldsymbol{\mu}_k \in \R^d$ are the component means, and
$\boldsymbol{\Sigma}_k \in \R^{d \times d}$ are positive definite
covariance matrices.
\end{definition}

\subsection{The EM algorithm}

The parameters $\theta = \{\pi_k, \boldsymbol{\mu}_k,
\boldsymbol{\Sigma}_k\}_{k=1}^K$ are estimated via the
\emph{Expectation-Maximisation} (EM) algorithm.

\textbf{E-step.}  Compute for each observation $\mathbf{x}_i$ the
\emph{responsibility} of component $k$:
\begin{equation}
    \gamma_{ik} = \mathbb{P}(Z_i = k \mid \mathbf{x}_i;\, \theta^{(t)})
    = \frac{\pi_k\,\mathcal{N}(\mathbf{x}_i;\,
      \boldsymbol{\mu}_k, \boldsymbol{\Sigma}_k)}
    {\sum_{j=1}^K \pi_j\,\mathcal{N}(\mathbf{x}_i;\,
      \boldsymbol{\mu}_j, \boldsymbol{\Sigma}_j)}.
    \label{eq:em-estep}
\end{equation}

\textbf{M-step.}  Update parameters via weighted averages:
\begin{align}
    N_k &= \sum_{i=1}^N \gamma_{ik}, \qquad
    \pi_k^{(t+1)} = \frac{N_k}{N}, \label{eq:em-pi}\\
    \boldsymbol{\mu}_k^{(t+1)} &= \frac{1}{N_k}
    \sum_{i=1}^N \gamma_{ik}\,\mathbf{x}_i, \label{eq:em-mu}\\
    \boldsymbol{\Sigma}_k^{(t+1)} &= \frac{1}{N_k}
    \sum_{i=1}^N \gamma_{ik}\,
    (\mathbf{x}_i - \boldsymbol{\mu}_k^{(t+1)})
    (\mathbf{x}_i - \boldsymbol{\mu}_k^{(t+1)})^\top.
    \label{eq:em-sigma}
\end{align}

\begin{remark}
K-means is a degenerate special case of EM for GMMs: setting all
covariance matrices to $\boldsymbol{\Sigma}_k = \varepsilon I$ and
forcing hard assignments ($\gamma_{ik} \in \{0,1\}$) recovers the
K-means update exactly.
\end{remark}

\subsection{Choosing $K$: BIC and AIC}

For GMMs the log-likelihood is monotone in $K$, so one uses penalised
criteria:
\begin{align}
    \mathrm{AIC}(K) &= -2\,\ell(\hat{\theta}_K) + 2\,p_K,
    \label{eq:aic}\\
    \mathrm{BIC}(K) &= -2\,\ell(\hat{\theta}_K) + p_K\,\log N,
    \label{eq:bic}
\end{align}
where $p_K = (K-1) + Kd + K\frac{d(d+1)}{2}$ is the number of free
parameters.  The BIC is preferred for clustering due to its consistency
as $N \to \infty$.

\section{Adaptive Simulation Allocation}
\label{sec:adaptive-allocation}

Let $\hat{\sigma}_k^2$ denote the within-cluster payoff variance
estimated from a pilot simulation of $N_0$ paths, and $p_k = N_k/N_0$
the empirical cluster weight.  The \emph{Neyman-optimal} allocation
assigns to cluster $k$:
\begin{equation}
    n_k^* = N \cdot \frac{p_k\,\hat{\sigma}_k}
    {\sum_{j=1}^K p_j\,\hat{\sigma}_j},
    \label{eq:neyman-cluster}
\end{equation}
yielding the adaptive estimator:
\begin{equation}
    \hat{V}_0^{ADAPT} = e^{-rT} \sum_{k=1}^{K} p_k
    \left(\frac{1}{n_k^*}\sum_{i \in C_k} \Phi(\omega_i)\right),
    \label{eq:adaptive-estimator}
\end{equation}
with variance:
\begin{equation}
    \Var\!\left[\hat{V}_0^{ADAPT}\right]
    = \frac{e^{-2rT}}{N}
    \left(\sum_{k=1}^{K} p_k\,\hat{\sigma}_k\right)^2.
    \label{eq:adaptive-variance}
\end{equation}

\section{Local Control Variates and Variance Reduction}
\label{sec:local-cv}

\subsection{Local control variates}

Rather than fitting a single global coefficient $c^*$, one fits a
separate coefficient within each cluster:
\begin{equation}
    c_k^* = \frac{\widehat{\Cov}_k[\Phi,\, Y]}
    {\widehat{\Var}_k[Y]},
    \label{eq:local-cv-coeff}
\end{equation}
yielding the local control variate estimator:
\begin{equation}
    \hat{V}_0^{LCV} = e^{-rT} \sum_{k=1}^{K} p_k \cdot
    \frac{1}{n_k} \sum_{i \in C_k}
    \bigl[\Phi(\omega_i) - c_k^*\,(Y_i - \mu_Y)\bigr].
    \label{eq:local-cv-estimator}
\end{equation}

\subsection{Why clustering amplifies variance reduction}

The \emph{law of total variance} decomposes payoff variance as:
\begin{equation}
    \Var[\Phi] = \underbrace{\mathbb{E}_k\bigl[\Var[\Phi \mid C_k]\bigr]}_{\text{within-cluster}}
    + \underbrace{\Var_k\bigl[\mathbb{E}[\Phi \mid C_k]\bigr]}_{\text{between-cluster}}.
    \label{eq:total-variance}
\end{equation}

Clustering attacks both terms simultaneously: stratification eliminates
the between-cluster component, while local control variates reduce the
within-cluster component by exploiting higher within-cluster
correlations.  The residual variance satisfies:
\[
    \sum_{k=1}^K p_k\,\sigma_k^2\,(1 - \rho_k^2)
    \;\leq\; \Var[\Phi]\,(1 - \rho_{\Phi,Y}^2),
\]
where $\rho_k = \mathrm{Corr}_k[\Phi, Y]$.

\section{Summary}
\label{sec:clustering-summary}

The complete pipeline consists of four steps: (1) feature extraction,
(2) standardisation, (3) clustering via K-means or GMM, and (4) adaptive
Neyman allocation combined with local control variates.  The variance
decomposition~\eqref{eq:total-variance} shows why these two mechanisms
act simultaneously and produce gains substantially greater than either
technique alone.

% ── Chapter 4 ───────────────────────────────────────────────
\chapter{Random Forests for Payoff Approximation and Control Variates}
\label{ch:random-forests}

Chapter~\ref{ch:clustering} showed that grouping simulated paths into
homogeneous clusters amplifies the effectiveness of control variates.
The remaining bottleneck is the choice of control $Y$ itself: in the
classical framework, $Y$ must be chosen by hand.  This chapter proposes
to \emph{learn} the control variate function directly from simulated
data using \emph{Random Forests} --- an ensemble method that is robust,
fast to train, and interpretable.

\section{Decision Trees}
\label{sec:decision-trees}

\subsection{Formal definition}

\begin{definition}[Regression tree]
A regression tree with $L$ leaves partitions $\R^d$ into $L$ disjoint
rectangular regions $R_1, \ldots, R_L$ and approximates the target
function by:
\begin{equation}
    \hat{f}(\mathbf{x}) = \sum_{\ell=1}^{L} c_\ell\,
    \mathbf{1}\{\mathbf{x} \in R_\ell\},
    \label{eq:tree}
\end{equation}
where $c_\ell = \frac{1}{|R_\ell|}\sum_{i:\,\mathbf{x}_i \in R_\ell}
\Phi_i$ is the within-leaf mean.
\end{definition}

The CART algorithm (Breiman et al., 1984) builds the tree by
recursively choosing the split $(\ell^*, s^*)$ that minimises the total
within-region residual sum of squares:
\begin{equation}
    \mathcal{L}(\ell, s) =
    \sum_{i:\,\mathbf{x}_i \in R^-} (\Phi_i - \bar{\Phi}^-)^2
    + \sum_{i:\,\mathbf{x}_i \in R^+} (\Phi_i - \bar{\Phi}^+)^2.
    \label{eq:cart-criterion}
\end{equation}

\subsection{Illustrative example}

To make the algorithm concrete, Table~\ref{tab:cart-example} presents a
pilot simulation of $N_0 = 6$ paths of a European call with $K = 100$,
described by the terminal value $S_T$ and the realised volatility
$\hat{\sigma}^2$.

% Table uses booktabs — consistent with the rest of the document
\begin{table}[H]
\centering
\caption{Pilot paths for the CART illustration.  The optimal first
         split $S_T \leq 100$ perfectly separates zero-payoff paths
         from positive-payoff paths.}
\label{tab:cart-example}
\renewcommand{\arraystretch}{1.3}
\begin{tabular}{cccr}
\toprule
\textbf{Path} $i$ & $S_T$ & $\hat{\sigma}^2$ & \textbf{Payoff} $(S_T - 100)^+$ \\
\midrule
1 & 112 & 0.06 & 12 \\
2 & 95  & 0.04 & 0  \\
3 & 108 & 0.08 & 8  \\
4 & 91  & 0.03 & 0  \\
5 & 103 & 0.05 & 3  \\
6 & 88  & 0.02 & 0  \\
\bottomrule
\end{tabular}
\end{table}

\section{Random Forests}
\label{sec:rf}

\begin{definition}[Random Forest]
A Random Forest with $B$ trees and $m$ candidate features per split is
the ensemble:
\begin{equation}
    \hat{f}^{RF}(\mathbf{x}) = \frac{1}{B}\sum_{b=1}^{B}
    \hat{f}_b^{m}(\mathbf{x}),
    \label{eq:rf}
\end{equation}
where $\hat{f}_b^m$ is a fully grown decision tree trained on the $b$-th
bootstrap sample, with $m$ randomly selected features considered at each
split.  The standard choice is $m = \lfloor d/3 \rfloor$ for regression.
\end{definition}

Random Forests enjoy three properties of particular relevance here.
\emph{Consistency}: the predictor $\hat{f}^{RF}(\mathbf{x})$ converges
to $\mathbb{E}[\Phi \mid \mathbf{x}]$ as $N, B \to \infty$ (Scornet et
al., 2015).  \emph{Out-of-bag error estimation}: roughly $37\%$ of
observations are unused per tree, providing a free validation estimate:
\begin{equation}
    \mathrm{OOB} = \frac{1}{N}\sum_{i=1}^{N}
    \left(\Phi_i - \hat{f}^{OOB}(\mathbf{x}_i)\right)^2.
    \label{eq:oob}
\end{equation}
\emph{Feature importance}: the mean decrease in impurity (MDI) ranks
features by their total contribution to splits across all trees.

\section{Random Forest Control Variates}
\label{sec:rf-cv}

The two-phase procedure runs as follows.  In the \emph{training phase},
one fits a Random Forest on a pilot sample $\{(\mathbf{x}_i,
\Phi_i)\}_{i=1}^{N_0}$ and estimates the control mean:
\begin{equation}
    \hat{\mu}_Y = \frac{1}{N_0}\sum_{i=1}^{N_0}
    \hat{f}^{OOB}(\mathbf{x}_i).
    \label{eq:rf-control-mean}
\end{equation}
In the \emph{pricing phase}, one evaluates the fitted forest on $N$
fresh paths:
\begin{equation}
    \hat{V}_0^{RF} = \frac{e^{-rT}}{N}\sum_{i=1}^{N}
    \left[\Phi(\omega_i) - c^*\,
    \bigl(\hat{f}^{RF}(\mathbf{x}_i) - \hat{\mu}_Y\bigr)\right],
    \label{eq:rf-cv-estimator}
\end{equation}
with variance reduction factor $(1 - \rho^2_{\Phi,\hat{f}^{RF}})$
approaching zero as the forest approximation improves.

The local variant combines clustering and RF control variates:
\begin{equation}
    \hat{V}_0^{LRF} = e^{-rT}\sum_{k=1}^{K} p_k \cdot
    \frac{1}{n_k}\sum_{i \in C_k}
    \left[\Phi(\omega_i) - c_k^*\,
    \bigl(\hat{f}_k^{RF}(\mathbf{x}_i) - \hat{\mu}_{Y,k}\bigr)\right].
    \label{eq:rf-local-estimator}
\end{equation}

\section{Comparison with Alternative Approaches}
\label{sec:rf-comparison}

Table~\ref{tab:method-comparison} summarises the three control variate
construction methods along the dimensions most relevant to a production
Monte Carlo engine.

\begin{table}[H]
\centering
\caption{Comparison of control variate construction methods.  Variance
         reduction ratings are qualitative assessments based on empirical
         results reported in Parts~II and~III.}
\label{tab:method-comparison}
\renewcommand{\arraystretch}{1.5}
\begin{tabular}{lcccc}
\toprule
\textbf{Method} & \textbf{Flexibility}
                & \textbf{Training cost}
                & \textbf{Interpretability}
                & \textbf{Variance reduction} \\
\midrule
Regression MC    & Low    & Very low  & High   & Low--Medium \\
Random Forest    & Medium & Low       & Medium & Medium--High \\
Neural Network   & High   & High      & Low    & High \\
\bottomrule
\end{tabular}
\end{table}

\section{Summary}
\label{sec:rf-summary}

This chapter developed Random Forests as a principled, data-driven method
for constructing control variates in Monte Carlo pricing.  The key
contributions are: (1) regression trees as adaptive stratification,
(2) ensemble averaging for variance reduction, (3) RF control variates
with convergence guarantee, (4) within-leaf variance for detecting
high-variance regions, and (5) local RF control variates that combine
clustering with data-driven control construction.

% ── Chapter 5 ───────────────────────────────────────────────
\chapter{Neural Networks for Variance Reduction}
\label{ch:neural-networks}

Chapter~\ref{ch:random-forests} showed that Random Forests can
approximate $\mathbb{E}^{\Q}[\Phi \mid \mathbf{x}]$ and use this
approximation as a control variate.  Their expressive power is
fundamentally limited, however, by the piecewise-constant nature of tree
predictions.  Neural networks offer a fundamentally different function
approximation paradigm: they compose simple non-linear transformations
in sequence, producing smooth, flexible approximations that can represent
virtually any continuous function --- a property known as \emph{universal
approximation}.

\section{From a Single Neuron to a Deep Network}
\label{sec:neuron}

\subsection{The artificial neuron}

The basic computational unit takes a vector of inputs
$\mathbf{x} \in \R^d$, computes a weighted sum, adds a bias, and passes
the result through a non-linear \emph{activation function}
$\sigma : \R \to \R$:
\begin{equation}
    z = \sigma\!\left(\mathbf{w}^\top \mathbf{x} + b\right),
    \label{eq:neuron}
\end{equation}
where $\mathbf{w} \in \R^d$ and $b \in \R$ are the learnable parameters.
Standard choices include: \emph{Sigmoid}
$\sigma(z) = (1 + e^{-z})^{-1}$;
\emph{Hyperbolic tangent} $\sigma(z) = \tanh(z)$;
\emph{ReLU} $\sigma(z) = \max(0,z)$, which is the most widely used in
modern deep networks.

\subsection{The feedforward network}

A \emph{feedforward neural network} (MLP) arranges neurons in layers.
The output of the $\ell$-th hidden layer is:
\begin{equation}
    \mathbf{h}^{(\ell)} = \sigma\!\left(W^{(\ell)}\,\mathbf{h}^{(\ell-1)}
    + \mathbf{b}^{(\ell)}\right), \qquad \ell = 1, \ldots, L-1,
    \label{eq:hidden-layer}
\end{equation}
with final output $g_\theta(\mathbf{x}) = W^{(L)}\,\mathbf{h}^{(L-1)}
+ \mathbf{b}^{(L)}$, where
$\theta = \{W^{(\ell)}, \mathbf{b}^{(\ell)}\}_{\ell=1}^{L}$ denotes
the full set of trainable parameters.

\section{Universal Approximation}
\label{sec:ua}

\begin{theorem}[Universal Approximation, Cybenko 1989]
Let $\sigma : \R \to \R$ be a continuous, bounded, non-constant
function.  For any continuous function $f : [0,1]^d \to \R$ and any
$\varepsilon > 0$, there exists a feedforward network $g_\theta$ with
one hidden layer such that:
\begin{equation}
    \sup_{\mathbf{x} \in [0,1]^d}
    \left|f(\mathbf{x}) - g_\theta(\mathbf{x})\right| < \varepsilon.
    \label{eq:uat}
\end{equation}
\end{theorem}

The theorem guarantees the \emph{existence} of such a network but says
nothing about whether training will \emph{find} it --- a distinction
with important practical implications (see Section~\ref{sec:nn-limits}).

\section{Training Neural Networks}
\label{sec:training}

\subsection{Loss functions}

For regression, the standard choice is the \emph{mean squared error}:
\begin{equation}
    \mathcal{L}_{MSE}(\theta) = \frac{1}{N_0}\sum_{i=1}^{N_0}
    \left(\Phi_i - g_\theta(\mathbf{x}_i)\right)^2.
    \label{eq:mse-loss}
\end{equation}
In Section~\ref{sec:nn-cv}, we use instead the \emph{variance loss}:
\begin{equation}
    \mathcal{L}_{VAR}(\theta, c) = \frac{1}{N_0}\sum_{i=1}^{N_0}
    \left(\Phi_i - c\,g_\theta(\mathbf{x}_i)\right)^2,
    \label{eq:var-loss}
\end{equation}
minimised jointly over $\theta$ and the scalar $c$.

\subsection{Backpropagation}

Define the error signal $\boldsymbol{\delta}^{(\ell)} =
\partial \mathcal{L} / \partial \mathbf{z}^{(\ell)}$.  The
backpropagation recursion is:
\begin{align}
    \boldsymbol{\delta}^{(L)} &= \frac{\partial \mathcal{L}}
    {\partial g_\theta}, \label{eq:bp-output}\\
    \boldsymbol{\delta}^{(\ell)} &= \left(W^{(\ell+1)}\right)^\top
    \boldsymbol{\delta}^{(\ell+1)} \odot \sigma'(\mathbf{z}^{(\ell)}),
    \label{eq:bp-hidden}\\
    \frac{\partial \mathcal{L}}{\partial W^{(\ell)}} &=
    \boldsymbol{\delta}^{(\ell)} \left(\mathbf{h}^{(\ell-1)}\right)^\top.
    \label{eq:bp-grad}
\end{align}

\subsection{Adam optimiser}

The Adam optimiser (Kingma \& Ba, 2015) adapts the learning rate per
parameter:
\begin{align}
    m_t &= \beta_1\,m_{t-1} + (1 - \beta_1)\,\nabla_\theta\mathcal{L}_t,
    \label{eq:adam-m}\\
    v_t &= \beta_2\,v_{t-1} + (1 - \beta_2)\,
    (\nabla_\theta\mathcal{L}_t)^2,
    \label{eq:adam-v}\\
    \theta^{(t+1)} &= \theta^{(t)} - \frac{\eta}{\sqrt{\hat{v}_t}
    + \varepsilon}\,\hat{m}_t,
    \label{eq:adam-update}
\end{align}
with default values $\beta_1 = 0.9$, $\beta_2 = 0.999$,
$\varepsilon = 10^{-8}$.

\section{Neural Network Control Variates}
\label{sec:nn-cv}

\subsection{The variance minimisation objective}

The goal is to find $g_\theta$ minimising the residual variance:
\begin{equation}
    \min_{\theta,\, c}\;\Var^{\Q}\!\left[\Phi(\omega)
    - c\,g_\theta(\mathbf{x}(\omega))\right],
    \label{eq:nn-cv-objective}
\end{equation}
which is equivalent to maximising the correlation
$\rho_{\Phi, g_\theta}$.  The optimal coefficient at fixed $\theta$ is:
\begin{equation}
    c^* = \frac{\sum_{i=1}^{N_0} \Phi_i\,g_\theta(\mathbf{x}_i)
    - N_0\,\bar{\Phi}\,\overline{g_\theta}}
    {\sum_{i=1}^{N_0} g_\theta(\mathbf{x}_i)^2
    - N_0\,\overline{g_\theta}^2}.
    \label{eq:nn-c-star}
\end{equation}

\subsection{Architecture}

Table~\ref{tab:nn-architecture} summarises the retained architecture
for European options under Heston with feature vector
$\tilde{\mathbf{x}} = (S_T, V_T, \bar{V}) \in \R^3$.

\begin{table}[H]
\centering
\caption{Neural network architecture for the control variate
         $g_\theta$.  Input: standardised feature vector
         $\tilde{\mathbf{x}} \in \R^d$ with $d = 3$.}
\label{tab:nn-architecture}
\renewcommand{\arraystretch}{1.4}
\begin{tabular}{lcccc}
\toprule
\textbf{Layer} & \textbf{Input} & \textbf{Output}
               & \textbf{Activation} & \textbf{Parameters} \\
\midrule
Input    & $d = 3$ & 3  & ---  & 0 \\
Hidden 1 & 3       & 64 & ReLU & $3 \times 64 + 64 = 256$ \\
Hidden 2 & 64      & 64 & ReLU & $64 \times 64 + 64 = 4{,}160$ \\
Hidden 3 & 64      & 64 & ReLU & $64 \times 64 + 64 = 4{,}160$ \\
Output   & 64      & 1  & None & $64 \times 1 + 1 = 65$ \\
\midrule
\textbf{Total} & & & & $\mathbf{8{,}641}$ \\
\bottomrule
\end{tabular}
\end{table}

\subsection{The training and pricing pipeline}

\begin{enumerate}
    \item \textbf{Pilot simulation}: simulate $N_0$ paths, compute
          features $\mathbf{x}_i$ and payoffs $\Phi_i$.
    \item \textbf{Standardisation}: standardise features to zero mean
          and unit variance.
    \item \textbf{Training}: minimise \eqref{eq:var-loss} via Adam,
          mini-batches of size $256$, $200$ epochs, $L^2$
          regularisation $\lambda = 10^{-4}$, dropout $p_{drop} = 0.1$.
    \item \textbf{Validation}: monitor out-of-sample variance reduction
          on a held-out subset.
    \item \textbf{Pricing}: simulate $N$ fresh paths and compute:
\begin{equation}
    \hat{V}_0^{NN} = \frac{e^{-rT}}{N}\sum_{i=1}^{N}
    \left[\Phi(\omega_i) - c^*\,
    \bigl(g_\theta(\mathbf{x}_i) - \hat{\mu}_{g}\bigr)\right].
    \label{eq:nn-cv-estimator}
\end{equation}
\end{enumerate}

\section{Application to High-Dimensional Finance}
\label{sec:nn-finance}

For a $d$-dimensional basket option with payoff
$(\bar{S}_T - K)^+ = \bigl(\frac{1}{d}\sum_{j=1}^d S_T^{(j)} - K\bigr)^+$
under multi-dimensional Black--Scholes:
\begin{equation}
    \dd S_t^{(j)} = r\,S_t^{(j)}\,\dd t
    + \sigma_j\,S_t^{(j)}\,\dd W_t^{(j)}, \qquad j = 1, \ldots, d,
    \label{eq:multi-bs}
\end{equation}
with $\langle W^{(j)}, W^{(k)}\rangle_t = \rho_{jk}\,t$, a three-layer
network of width $128$ trained on $N_0 = 10^5$ paths typically achieves
$\rho > 0.95$ for $d = 5$ to $20$, reducing variance by a factor
exceeding $20\times$.

\section{Practical Limitations}
\label{sec:nn-limits}

Neural network control variates face five practical challenges that
practitioners must weigh carefully.  \textbf{Overfitting} can be
detected via held-out validation and mitigated by $L^2$ regularisation
and dropout.  \textbf{Training instability} arises for discontinuous
payoffs (digitals, barriers), requiring smooth indicator approximations.
\textbf{Computational cost} is on the order of $10$--$60$ seconds on a
modern CPU for the architectures above, reduced by $10$--$100\times$
with GPU acceleration.  \textbf{Interpretability} is limited relative
to Random Forests; SHAP values (Lundberg \& Lee, 2017) provide post-hoc
attribution at additional cost.  \textbf{No-arbitrage constraints} can
be enforced via a Lagrange multiplier on the martingale condition
$\mathbb{E}^{\Q}[g_\theta(\mathbf{x})] = 0$.

\section{Summary}
\label{sec:nn-summary}

This chapter developed neural networks as the most powerful and flexible
control variate construction method in our hierarchy, deriving the
variance minimisation objective, the backpropagation training procedure,
and the full pricing pipeline.  Together, Chapters~\ref{ch:clustering}
to~\ref{ch:neural-networks} constitute the complete machine learning
pipeline:
\begin{enumerate}
    \item Feature extraction and normalisation
    \item Path clustering via K-means or GMM
    \item Adaptive Neyman allocation
    \item Random Forest local control variates
    \item Neural network global control variates
\end{enumerate}
Part~III implements and benchmarks all five steps against the classical
methods of Part~I across a range of option types and underlying models.

% ============================================================
%  APPENDIX
%  The \appendix command switches chapter labels to A, B, C ...
%  \appendixpage inserts a divider page ("Appendices").
%  \input pulls the standalone appendix.tex without a page break.
% ============================================================
\appendix
\appendixpage
\addappheadtotoc

% ============================================================
%  APPENDIX
%  Stratified Monte Carlo Pricing and Variance Reduction
%  via K-Means Clustering and Machine Learning Control Variates
%
%  Benjamin Sulpice & Enzo Thuilliez — May 2026
%  Supervised by Professor Arnaud DUFAYS
% ============================================================
%
%  This appendix collects the four publication-quality figures
%  generated by the Python pipeline (main.py) and the critical
%  code extracts underlying the three machine-learning steps of
%  the framework: K-Means/Neyman allocation (Step 3), local
%  Random Forest control variates (Step 4), and the global
%  neural-network variance-loss optimiser (Step 5).
%
%  All figures are formatted consistently with the main text
%  (monochrome/grayscale, serif typeface, minimalist layout).
%  All code blocks use the \texttt{listings} environment with
%  monochrome syntax highlighting: bold keywords, italic
%  comments, no colour.
% ============================================================

% ── Appendix counter setup ───────────────────────────────────
% Chapter, figure, table, and equation counters are reset per
% appendix chapter and labelled A.1, B.1, etc. automatically
% when this file is included after \appendix in main.tex.

% ============================================================
\chapter{Numerical Figures}
\label{app:figures}
% ============================================================

The figures presented in this appendix constitute the four primary
visual outputs of the five-step pipeline implemented in \texttt{main.py}.
Each figure is generated under identical model parameters
($S_0 = K = 100$, $r = 0.05$, $\sigma = 0.20$, $T = 1$, $M = 252$
daily time steps) and saved in publication-quality grayscale at
200\,dpi.  The sequence of figures mirrors the logical progression
of the pipeline: from raw path simulation (Figure~\ref{fig:trajectoires}),
through empirical convergence benchmarks (Figure~\ref{fig:rmse}),
to the diagnostic outputs of the clustering stage
(Figure~\ref{fig:clusters}) and the neural-network optimiser
(Figure~\ref{fig:nn_loss}).

% ------------------------------------------------------------
\section{Figure 1: Simulated Asset Trajectories}
\label{app:fig1}
% ------------------------------------------------------------

\paragraph{Context.}
The first step of the pipeline generates $N$ Monte Carlo paths
of the underlying asset under the risk-neutral Geometric Brownian
Motion (GBM) dynamics:
\[
  S_t = S_0 \exp\!\left[\left(r - \tfrac{1}{2}\sigma^2\right)t
  + \sigma W_t^{\mathbb{Q}}\right], \quad t \in [0, T],
\]
where $W^{\mathbb{Q}}$ is a standard Brownian motion under the
pricing measure $\mathbb{Q}$.  Figure~\ref{fig:trajectoires}
displays a representative sample of $N = 200$ such paths,
colour-coded by their terminal moneyness status.  This
visualisation serves two purposes: (i) it validates the
correctness of the discretisation scheme and the resulting
distributional properties of $S_T$, and (ii) it provides an
immediate geometric illustration of the heterogeneity that
motivates trajectory clustering --- paths that expire deeply
in the money exhibit systematically different path
characteristics (higher mean, lower realised variance) than
out-of-the-money paths, justifying the partition of the
simulation space into homogeneous strata.

\begin{figure}[H]
  \centering
  \includegraphics[width=\linewidth]{figure_1_trajectoires.png}
  \caption{%
    \textbf{Simulated GBM trajectories categorised by terminal
    moneyness.}
    Two hundred paths of Geometric Brownian Motion are displayed
    under parameters $S_0 = K = 100$, $r = 0.05$,
    $\sigma = 0.20$, $T = 1$, $M = 252$.
    \textit{Left panel}: European call perspective.
    Dark paths expire in-the-money ($S_T > K$); light paths
    expire out-of-the-money ($S_T \leq K$).  The dashed
    horizontal line marks the strike $K$.
    \textit{Right panel}: European put perspective, with dark
    paths now corresponding to $S_T < K$.
    The bimodal structure of the terminal distribution under
    the moneyness partition provides the geometric motivation
    for K-Means clustering in the five-dimensional feature
    space $(S_T,\,\bar{S},\,\widehat{\sigma}^2,\,S_{\max},\,S_{\min})$.%
  }
  \label{fig:trajectoires}
\end{figure}

% ------------------------------------------------------------
\section{Figure 2: Empirical RMSE Convergence Rates}
\label{app:fig2}
% ------------------------------------------------------------

\paragraph{Context.}
The fundamental limitation of plain Monte Carlo pricing is its
slow convergence rate: the Root Mean Squared Error (RMSE) of the
estimator $\hat{V}_0^{MC} = N^{-1}\sum_{i=1}^N \phi_i$ decreases
as $\mathcal{O}(N^{-1/2})$, where $N$ is the number of simulated
paths.  Every variance reduction technique in the pipeline can be
understood as an attempt to reduce the constant $C$ in
$\mathrm{RMSE} \approx C \cdot N^{-1/2}$, or equivalently, to
reduce the variance $\mathrm{Var}(\phi_i)$ of the effective payoff
being averaged.

Figure~\ref{fig:rmse} benchmarks five architectural variants of
the Monte Carlo estimator across sample sizes
$N \in \{500,\,1{,}000,\,2{,}000,\,5{,}000,\,10{,}000,\,20{,}000\}$.
For each $(N, \text{method})$ pair, $n = 100$ independent
replications are performed, and the RMSE is estimated as
$\sqrt{n^{-1}\sum_{j=1}^n (\hat{V}_{0,j} - V_0^*)^2}$, where
$V_0^* = V_0^{BS}$ is the analytical Black--Scholes price.  The
log-log scale allows the empirical convergence slopes to be
compared directly against the $\mathcal{O}(N^{-1/2})$ reference
line (shown as a dashed gray line).  A steeper slope indicates
super-$\sqrt{N}$ convergence in the effective sample size, which
is equivalent to a variance reduction factor that grows with $N$.

\begin{figure}[H]
  \centering
  \includegraphics[width=0.88\linewidth]{figure_2_RMSE.png}
  \caption{%
    \textbf{Log-log empirical RMSE of five Monte Carlo
    estimators across sample sizes $N \in [500,\,20{,}000]$.}
    Each RMSE value is computed over $n = 100$ independent
    replications.  Estimators shown: Plain MC (solid, circles),
    Antithetic Variates (dashed, squares), Stratified K-Means
    (dash-dot, triangles), Random Forest control variate
    (dotted, diamonds), and Neural Network control variate
    (dash-dot-dot, stars).  The gray dashed line represents
    the theoretical $\mathcal{O}(N^{-1/2})$ reference slope.
    Machine-learning control variates (RF, NN) achieve a
    downward vertical shift of the RMSE curve, corresponding
    to a multiplicative variance reduction factor that is
    independent of $N$ and preserved uniformly across the
    sample-size range tested.  The stratified estimator
    exhibits a reduction consistent with the theoretical
    within-stratum variance decomposition.%
  }
  \label{fig:rmse}
\end{figure}

% ------------------------------------------------------------
\section{Figure 3: K-Means Cluster Structure in Feature Space}
\label{app:fig3}
% ------------------------------------------------------------

\paragraph{Context.}
The K-Means algorithm (Step 3 of the pipeline) partitions the
$N_{\text{pilot}}$ simulated paths into $K$ clusters by minimising
the total within-cluster variance of the normalised feature
vectors $\mathbf{x}_i \in \mathbb{R}^5$.  The features are:
terminal value $S_T$, arithmetic path mean $\bar{S}$, realised
variance $\widehat{\sigma}^2$, path maximum $S_{\max}$, and path
minimum $S_{\min}$.  Although the clustering is performed in this
full five-dimensional space, Figure~\ref{fig:clusters} projects
the assignments onto the $(S_T, \bar{S})$ plane, which captures
the dominant sources of payoff heterogeneity for standard European
and path-dependent options.

The figure demonstrates that K-Means recovers economically
meaningful regimes: clusters in the left region of the plot
(low $S_T$, low $\bar{S}$) correspond to deeply out-of-the-money
paths with near-zero call payoff and high put payoff; clusters in
the right region correspond to in-the-money paths; and clusters
straddling the strike line $K = 100$ (dashed) capture the
near-money paths, which contribute the highest within-cluster
payoff variance and therefore receive the largest Neyman allocation.

\begin{figure}[H]
  \centering
  \includegraphics[width=0.85\linewidth]{figure_3_clusters.png}
  \caption{%
    \textbf{K-Means cluster assignments projected onto the
    $(S_T,\,\bar{S})$ feature plane.}
    $N_{\text{pilot}} = 20{,}000$ GBM paths are clustered into
    $K = 8$ strata using K-Means++ initialisation.  Each point
    represents one simulated path, plotted at its terminal value
    $S_T$ (horizontal axis) against its arithmetic mean $\bar{S}$
    (vertical axis).  Distinct grayscale shades and marker
    symbols identify the eight cluster assignments.  The dashed
    vertical line marks the strike $K = 100$.  The cluster
    boundaries reveal a structure aligned with the moneyness
    partition and with the path-average distribution, confirming
    that the unsupervised algorithm recovers economically
    interpretable regimes without any explicit payoff information.
    The Neyman optimal allocation $n_k^* \propto p_k\,\sigma_k$
    subsequently concentrates the simulation budget on
    near-money clusters, where within-cluster payoff variance
    $\sigma_k$ is highest.%
  }
  \label{fig:clusters}
\end{figure}

% ------------------------------------------------------------
\section{Figure 4: Neural Network Variance-Loss Trajectory}
\label{app:fig4}
% ------------------------------------------------------------

\paragraph{Context.}
Step 5 of the pipeline trains a feedforward neural network
$g_\theta : \mathbb{R}^5 \to \mathbb{R}$ to serve as a global
control variate.  The network is optimised with respect to the
\emph{custom variance loss}:
\[
  \mathcal{L}(\theta, c)
  \;=\;
  \mathbb{E}\!\left[\bigl(\phi - c\,g_\theta(\mathbf{x})\bigr)^2\right]
  \;=\;
  \Var(\phi) - 2c\,\Cov(\phi, g_\theta) + c^2\,\Var(g_\theta),
\]
where $\phi$ is the (discounted) path payoff and $\mathbf{x}$
is the corresponding feature vector.  The scalar coefficient
$c$ is updated in closed form after each epoch:
\[
  c^*_{\text{epoch}}
  = \frac{\widehat{\Cov}(\phi,\,g_\theta)}{\widehat{\Var}(g_\theta)},
\]
which is precisely the optimal OLS coefficient minimising
$\mathcal{L}(\theta, c)$ with respect to $c$ for fixed $\theta$.
This alternating optimisation ensures that the loss converges
monotonically to the residual variance $\Var(\phi - c^* g_\theta)$,
which equals $(1 - \rho^2)\,\Var(\phi)$ at the global optimum,
where $\rho = \mathrm{Corr}(\phi, g_\theta)$.

Figure~\ref{fig:nn_loss} traces this minimisation across
$200$ training epochs on $N_{\text{pilot}} = 20{,}000$ paths,
providing empirical confirmation that the custom loss
converges reliably under the Adam optimiser.

\begin{figure}[H]
  \centering
  \includegraphics[width=0.82\linewidth]{figure_4_nn_loss.png}
  \caption{%
    \textbf{Optimisation trajectory of the neural-network
    variance loss $\mathcal{L}(\theta, c)$ over $200$ epochs.}
    The network $g_\theta$ is a three-hidden-layer feedforward
    architecture with 64 units per layer and ReLU activations,
    trained via Adam ($\alpha = 10^{-3}$, $\lambda = 10^{-4}$
    weight decay, $p = 0.10$ dropout).  The solid black line
    plots the per-epoch batch-average loss; the dashed gray
    line shows the 10-epoch moving average.  The annotated
    terminal value confirms convergence to a stable minimum.
    The final loss corresponds to $\mathcal{L}^* = (1 -
    \rho^2)\,\widehat{\Var}(\phi)$, where $\rho > 0.90$ is the
    empirical correlation between payoffs and network outputs
    on the pilot set, implying a variance reduction factor
    exceeding $5\times$ relative to plain Monte Carlo.%
  }
  \label{fig:nn_loss}
\end{figure}

% ============================================================
\chapter{Python Code Extracts}
\label{app:code}
% ============================================================

This appendix chapter reproduces the three critical algorithmic
routines from \texttt{main.py} that implement Steps 3--5 of the
pipeline.  The extracts are presented in full, without
abbreviation, to allow complete reproducibility of the numerical
results reported in the main text.  Syntax highlighting is
restricted to bold keywords and italic comments, in keeping with
the monochrome typographic standard of this document.

All code is written in Python 3.10+ and requires the following
libraries: \texttt{numpy}, \texttt{scikit-learn},
\texttt{torch} (PyTorch), and \texttt{scipy}.  The functions
below are self-contained and callable once the upstream feature
matrix \texttt{X\_scaled} and pilot payoff vector \texttt{payoffs}
have been constructed (Steps 1--2).

% ------------------------------------------------------------
\section{K-Means Clustering and Neyman Allocation (Step 3)}
\label{app:code-kmeans}
% ------------------------------------------------------------

\paragraph{Mathematical background.}
Given a normalised feature matrix
$\mathbf{X} = \{\mathbf{x}_i\}_{i=1}^{N} \subset \mathbb{R}^5$,
K-Means++ solves the combinatorial optimisation problem:
\[
  \min_{\{C_k\}_{k=1}^K}
  \sum_{k=1}^K \sum_{\mathbf{x} \in C_k}
  \|\mathbf{x} - \boldsymbol{\mu}_k\|_2^2,
\]
where $\boldsymbol{\mu}_k = |C_k|^{-1}\sum_{\mathbf{x} \in C_k}\mathbf{x}$
is the centroid of cluster $k$.  Once the pilot paths are assigned
to clusters, the Neyman optimal budget allocation concentrates
simulation effort on the strata with the highest variance--weight
product:
\[
  n_k^* = N \cdot
  \frac{p_k\,\sigma_k}{\sum_{j=1}^K p_j\,\sigma_j},
  \qquad
  p_k = \frac{|C_k|}{N}, \quad
  \sigma_k = \widehat{\mathrm{Std}}(\phi \mid C_k).
\]
The following listing implements both the clustering and the
allocation calculation in a single coherent routine.

\begin{lstlisting}[
  language=Python,
  caption={K-Means++ clustering and Neyman optimal allocation
           (from \texttt{main.py}, Section~5).},
  label={lst:kmeans}
]
from sklearn.cluster import KMeans
import numpy as np

def cluster_paths_kmeans(X_scaled, K, random_state=42):
    """
    K-Means++ clustering on normalised feature matrix.

    Parameters
    ----------
    X_scaled     : (N, d) float ndarray  -- standardised features
    K            : int                   -- number of clusters
    random_state : int

    Returns
    -------
    labels   : (N,) int ndarray   -- cluster assignment per path
    centers  : (K, d) float ndarray -- centroid coordinates
    km_model : fitted KMeans object
    """
    km = KMeans(
        n_clusters=K,
        init="k-means++",   # smart initialisation (Arthur & Vassilvitskii, 2007)
        n_init=10,           # 10 restarts to avoid local minima
        random_state=random_state
    )
    labels = km.fit_predict(X_scaled)
    return labels, km.cluster_centers_, km


def neyman_allocation(payoffs, labels, K, N_total):
    """
    Compute Neyman-optimal simulation budget per cluster.

    Formula: n_k* = N * (p_k * sigma_k) / sum_j(p_j * sigma_j)

    Parameters
    ----------
    payoffs  : (N_pilot,) float ndarray -- pilot discounted payoffs
    labels   : (N_pilot,) int ndarray   -- cluster assignments
    K        : int                      -- number of clusters
    N_total  : int                      -- total pricing budget

    Returns
    -------
    allocation : (K,) int ndarray   -- paths allocated to each cluster
    p_k        : (K,) float ndarray -- empirical cluster weights
    sigma_k    : (K,) float ndarray -- within-cluster payoff std devs
    """
    N_pilot = len(payoffs)
    p_k     = np.zeros(K)
    sigma_k = np.zeros(K)

    # Step 3a: compute within-cluster statistics
    for k in range(K):
        mask   = labels == k
        n_k    = mask.sum()
        p_k[k] = n_k / N_pilot
        # Use ddof=1 (unbiased) standard deviation
        sigma_k[k] = payoffs[mask].std(ddof=1) if n_k >= 2 else 0.0

    # Step 3b: Neyman weights
    denom   = (p_k * sigma_k).sum()
    weights = (p_k * sigma_k / denom) if denom > 0 else p_k

    # Step 3c: round to integer counts, then correct rounding error
    allocation = np.maximum(
        np.round(weights * N_total).astype(int), 1
    )
    # Ensure allocation sums to exactly N_total
    allocation[np.argmax(weights)] += N_total - allocation.sum()

    return allocation, p_k, sigma_k
\end{lstlisting}

% ------------------------------------------------------------
\section{Local Random Forest Control Variate (Step 4)}
\label{app:code-rf}
% ------------------------------------------------------------

\paragraph{Mathematical background.}
For each cluster $k$, a Random Forest regressor
$\hat{f}_k : \mathbb{R}^5 \to \mathbb{R}$ is fitted on the
pilot paths belonging to $C_k$.  The locally corrected payoff
estimator for cluster $k$ is:
\[
  \tilde{\phi}_i^{(k)}
  = \phi_i - c_k^*\bigl(\hat{f}_k(\mathbf{x}_i) - \mu_{Y,k}\bigr),
  \qquad
  c_k^* = \frac{\widehat{\Cov}(\phi,\,\hat{f}_k)}
               {\widehat{\Var}(\hat{f}_k)},
  \quad
  \mu_{Y,k} = \mathbb{E}[\hat{f}_k(X) \mid X \in C_k],
\]
where $c_k^*$ is the OLS-optimal scalar coefficient and
$\mu_{Y,k}$ ensures that the control variate correction is
mean-zero, preserving the unbiasedness of the estimator.  The
global stratified price estimate combines the within-cluster
corrected means weighted by the empirical cluster proportions $p_k$.

\begin{lstlisting}[
  language=Python,
  caption={Local Random Forest control variate training and
           application (from \texttt{main.py}, Section~7).},
  label={lst:rf}
]
from sklearn.ensemble import RandomForestRegressor
import numpy as np

def rf_control_variate(
    X_pilot, payoffs_pilot, labels_pilot,
    X_pricing, labels_pricing, p_k, K,
    n_estimators=200, max_depth=None, random_state=42
):
    """
    Fit one Random Forest per cluster; use predictions as
    local control variates to reduce within-cluster variance.

    Parameters
    ----------
    X_pilot         : (N_pilot, d)   normalised pilot features
    payoffs_pilot   : (N_pilot,)     discounted pilot payoffs
    labels_pilot    : (N_pilot,)     cluster labels (pilot)
    X_pricing       : (N_price, d)   normalised pricing features
    labels_pricing  : (N_price,)     cluster labels (pricing)
    p_k             : (K,)           empirical cluster weights
    K               : int            number of clusters

    Returns
    -------
    estimate          : float   -- corrected price estimate
    residual_var_ratio: float   -- Var(corrected) / Var(raw)
    """
    forests    = {}   # one RF model per cluster
    ctrl_means = {}   # E[f_k(X) | X in C_k]  (for mean-zero correction)

    # ── Training phase: fit one RF per cluster ──────────────
    for k in range(K):
        mask = labels_pilot == k
        if mask.sum() < 10:           # insufficient data: skip
            forests[k]    = None
            ctrl_means[k] = 0.0
            continue

        rf = RandomForestRegressor(
            n_estimators=n_estimators,
            max_depth=max_depth,
            min_samples_leaf=5,       # prevents overfitting on small leaves
            n_jobs=-1,
            random_state=random_state
        )
        rf.fit(X_pilot[mask], payoffs_pilot[mask])
        forests[k]    = rf
        # Control variate mean (estimated on pilot training data)
        ctrl_means[k] = rf.predict(X_pilot[mask]).mean()

    # ── Pricing phase: apply local OLS correction ───────────
    corrected, raw = [], []

    for k in range(K):
        mask_tr = labels_pilot  == k
        mask_pr = labels_pricing == k
        n_tr    = mask_tr.sum()
        n_pr    = mask_pr.sum()

        if n_pr == 0:
            continue

        # Align pilot sample to pricing cluster size
        phi_pr = (payoffs_pilot[mask_tr][:n_pr]
                  if n_tr >= n_pr
                  else payoffs_pilot[mask_tr])
        if len(phi_pr) == 0:
            continue

        raw.extend(phi_pr.tolist())

        if forests[k] is None:        # fallback: no correction
            corrected.extend(phi_pr.tolist())
            continue

        X_pr   = X_pricing[mask_pr][:len(phi_pr)]
        Y_pr   = forests[k].predict(X_pr)
        mu_Y_k = ctrl_means[k]

        # OLS optimal coefficient: c* = Cov(phi, Y) / Var(Y)
        cov_kk = np.cov(phi_pr, Y_pr)[0, 1]
        var_Y  = Y_pr.var(ddof=1)
        c_k    = cov_kk / var_Y if var_Y > 1e-10 else 0.0

        # Apply mean-zero control variate correction
        corrected.extend(
            (phi_pr - c_k * (Y_pr - mu_Y_k)).tolist()
        )

    corrected = np.array(corrected)
    raw       = np.array(raw)
    raw_var   = raw.var(ddof=1) if raw.var(ddof=1) > 0 else 1.0

    return corrected.mean(), corrected.var(ddof=1) / raw_var
\end{lstlisting}

% ------------------------------------------------------------
\section{Neural Network Variance-Loss Optimiser (Step 5)}
\label{app:code-nn}
% ------------------------------------------------------------

\paragraph{Mathematical background.}
The global neural network control variate solves:
\[
  \min_{\theta \in \Theta,\, c \in \mathbb{R}}
  \mathcal{L}(\theta, c)
  = \mathbb{E}\!\left[
      \bigl(\phi - c\,g_\theta(\mathbf{x})\bigr)^2
    \right],
\]
which equals $\Var(\phi) - 2c\,\Cov(\phi, g_\theta) + c^2\,\Var(g_\theta)$.
The minimiser over $c$ at fixed $\theta$ is the OLS coefficient
$c^* = \Cov(\phi, g_\theta)/\Var(g_\theta)$, and at the global optimum
$\mathcal{L}^* = (1 - \rho^2)\Var(\phi)$ where $\rho = \mathrm{Corr}(\phi, g_\theta)$.
The alternating update --- gradient step in $\theta$, closed-form update
of $c$ --- ensures monotone decrease of the loss.  The network
architecture is:
\[
  g_\theta : \mathbb{R}^5
  \xrightarrow{\text{Linear}_{64}} \text{ReLU}
  \xrightarrow{\text{Linear}_{64}} \text{ReLU}
  \xrightarrow{\text{Linear}_{64}} \text{ReLU}
  \xrightarrow{\text{Linear}_{1}} \mathbb{R}.
\]
Weights are initialised with He (Kaiming) normal initialisation
for ReLU activations.

\begin{lstlisting}[
  language=Python,
  caption={Neural network architecture and custom variance-loss
           training routine (from \texttt{main.py}, Section~8).},
  label={lst:nn}
]
import torch
import torch.nn as nn
import torch.optim as optim
import numpy as np


class ControlVariateNet(nn.Module):
    """
    Feedforward network  g_theta : R^d -> R.

    Architecture  : Input -> [Linear(64) -> ReLU] x 3 -> Linear(1)
    Initialisation: He (Kaiming) normal for ReLU layers.
    """
    def __init__(self, d_in: int):
        super().__init__()
        self.net = nn.Sequential(
            nn.Linear(d_in, 64), nn.ReLU(),
            nn.Linear(64, 64),   nn.ReLU(),
            nn.Linear(64, 64),   nn.ReLU(),
            nn.Linear(64, 1),
        )
        # He initialisation: preserves variance across ReLU layers
        for m in self.modules():
            if isinstance(m, nn.Linear):
                nn.init.kaiming_normal_(m.weight, nonlinearity="relu")
                nn.init.zeros_(m.bias)

    def forward(self, x):
        return self.net(x).squeeze(-1)


def train_nn_control_variate(
    X_train, phi_train,
    n_epochs=200, batch_size=256, lr=1e-3,
    weight_decay=1e-4, dropout_p=0.1, device="cpu",
    verbose=False, record_loss=False,
):
    """
    Train ControlVariateNet with the custom variance loss
    L(theta, c) = E[(phi - c * g_theta(x))^2].

    The scalar c is updated in closed form after each epoch:
        c* = Cov(phi, g_theta) / Var(g_theta)

    Parameters
    ----------
    X_train      : (N, d) float ndarray  -- normalised features
    phi_train    : (N,)   float ndarray  -- discounted payoffs
    n_epochs     : int    training epochs
    batch_size   : int    mini-batch size (SGD)
    lr           : float  Adam learning rate
    weight_decay : float  L2 regularisation coefficient
    dropout_p    : float  hidden-layer dropout probability
    device       : str    'cpu' or 'cuda'
    record_loss  : bool   if True, return per-epoch loss history

    Returns
    -------
    net       : trained ControlVariateNet
    c_star    : float   -- optimal scalar coefficient
    mu_g      : float   -- E[g_theta(x)] on training set
    loss_hist : list    -- per-epoch losses (if record_loss=True)
    """
    X_t   = torch.tensor(X_train,  dtype=torch.float32).to(device)
    phi_t = torch.tensor(phi_train, dtype=torch.float32).to(device)
    N     = X_train.shape[0]

    net       = ControlVariateNet(X_train.shape[1]).to(device)
    optimiser = optim.Adam(
        net.parameters(), lr=lr, weight_decay=weight_decay
    )
    c_star    = 1.0          # initialise coefficient
    loss_hist = []

    for epoch in range(n_epochs):
        net.train()
        perm       = torch.randperm(N)
        epoch_loss = 0.0
        n_batches  = 0

        for start in range(0, N, batch_size):
            idx   = perm[start: start + batch_size]
            x_b   = X_t[idx]
            phi_b = phi_t[idx]

            optimiser.zero_grad()
            g_b = net(x_b)

            # Custom variance loss: Var(phi - c * g_theta)
            loss = ((phi_b - c_star * g_b) ** 2).mean()
            loss.backward()      # backpropagate through g_theta only
            optimiser.step()
            epoch_loss += loss.item()
            n_batches  += 1

        # Closed-form update of c after each full epoch
        net.eval()
        with torch.no_grad():
            g_all    = net(X_t)
            phi_mean = phi_t.mean()
            g_mean   = g_all.mean()
            # c* = Cov(phi, g) / Var(g)
            num = ((phi_t - phi_mean) * (g_all - g_mean)).mean().item()
            den = ((g_all - g_mean) ** 2).mean().item()
            if abs(den) > 1e-10:
                c_star = num / den

        avg_loss = epoch_loss / max(n_batches, 1)
        loss_hist.append(avg_loss)

        if verbose and (epoch + 1) % 50 == 0:
            print(
                f"  Epoch {epoch+1:3d}/{n_epochs}  "
                f"loss={avg_loss:.6f}  c*={c_star:.4f}"
            )

    # Final mean of g_theta (for unbiased price correction)
    net.eval()
    with torch.no_grad():
        mu_g = net(X_t).mean().item()

    return (net, c_star, mu_g, loss_hist) if record_loss \
           else (net, c_star, mu_g)
\end{lstlisting}

% ============================================================
\chapter{Supplementary Parameter Tables}
\label{app:params}
% ============================================================

Table~\ref{tab:params-model} summarises the model and simulation
parameters used throughout the numerical study.
Table~\ref{tab:params-nn} collects the hyperparameters of the
neural network control variate.

\begin{table}[H]
  \centering
  \caption{%
    \textbf{Model and simulation parameters.}
    Parameters are held fixed across all experiments unless
    otherwise noted in the main text.%
  }
  \label{tab:params-model}
  \begin{tabular}{llll}
    \toprule
    Symbol & Description & Value & Unit \\
    \midrule
    $S_0$ & Initial asset price & 100 & \$ \\
    $K$   & Strike price        & 100 & \$ \\
    $r$   & Risk-free rate      & 0.05 & ann.\ rate \\
    $\sigma$ & GBM volatility   & 0.20 & ann.\ vol \\
    $T$   & Maturity            & 1.0  & years \\
    $M$   & Time steps          & 252  & steps (daily) \\
    $N_{\text{pilot}}$  & Pilot paths & 20\,000 & paths \\
    $N_{\text{price}}$  & Pricing paths & 20\,000 & paths \\
    $K_{\text{clust}}$ & Number of clusters & 8 & --- \\
    \midrule
    $\kappa$ & Heston mean-reversion speed & 2.0  & --- \\
    $\theta$ & Heston long-run variance    & 0.04 & --- \\
    $\xi$    & Heston vol-of-vol           & 0.30 & --- \\
    $\rho$   & Heston correlation          & $-0.70$ & --- \\
    $V_0$    & Heston initial variance     & 0.04 & --- \\
    \bottomrule
  \end{tabular}
\end{table}

\begin{table}[H]
  \centering
  \caption{%
    \textbf{Neural network hyperparameters.}
    The architecture and training configuration are held fixed
    across all option types and model dynamics.%
  }
  \label{tab:params-nn}
  \begin{tabular}{lll}
    \toprule
    Hyperparameter & Description & Value \\
    \midrule
    Architecture   & Hidden layers $\times$ units & $3 \times 64$ \\
    Activation     & Hidden activation             & ReLU \\
    Initialisation & Weight initialisation         & He (Kaiming) normal \\
    Optimiser      & Gradient method               & Adam \\
    Learning rate  & $\alpha$                      & $10^{-3}$ \\
    Weight decay   & L2 regularisation $\lambda$   & $10^{-4}$ \\
    Dropout        & Hidden-layer dropout $p$      & $0.10$ \\
    Batch size     & Mini-batch size               & 256 \\
    Epochs         & Training epochs               & 200 \\
    Pilot size     & Pre-training paths $N_0$      & 50\,000 \\
    \bottomrule
  \end{tabular}
\end{table}

% ============================================================
%  BIBLIOGRAPHY
%  Placed after the appendix so the TOC sequence reads:
%  Chapters → Appendix → Bibliography.
%  Using the self-contained thebibliography environment.
% ============================================================
\cleardoublepage
\addcontentsline{toc}{chapter}{Bibliography}

\begin{thebibliography}{99}

\bibitem{Andersen2008}
Andersen, L.\ (2008).
\textit{Simple and efficient simulation of the Heston stochastic
  volatility model}.
Journal of Computational Finance, 11(3), 1--42.

\bibitem{ArthurVassilvitskii2007}
Arthur, D.\ \& Vassilvitskii, S.\ (2007).
\textit{k-means++: The advantages of careful seeding}.
In Proceedings of the 18th Annual ACM-SIAM Symposium on Discrete
Algorithms (SODA), 1027--1035.

\bibitem{BlackScholes1973}
Black, F.\ \& Scholes, M.\ (1973).
\textit{The pricing of options and corporate liabilities}.
Journal of Political Economy, 81(3), 637--654.

\bibitem{BreimanEtAl1984}
Breiman, L., Friedman, J.\ H., Olshen, R.\ A., \& Stone, C.\ J.\
(1984). \textit{Classification and Regression Trees}.
Wadsworth, Belmont, CA.

\bibitem{Breiman1996}
Breiman, L.\ (1996).
\textit{Bagging predictors}.
Machine Learning, 24(2), 123--140.

\bibitem{Breiman2001}
Breiman, L.\ (2001).
\textit{Random forests}.
Machine Learning, 45(1), 5--32.

\bibitem{Cybenko1989}
Cybenko, G.\ (1989).
\textit{Approximation by superpositions of a sigmoidal function}.
Mathematics of Control, Signals, and Systems, 2(4), 303--314.

\bibitem{DelbaenSchachermayer1994}
Delbaen, F.\ \& Schachermayer, W.\ (1994).
\textit{A general version of the fundamental theorem of asset pricing}.
Mathematische Annalen, 300(1), 463--520.

\bibitem{Glasserman2004}
Glasserman, P.\ (2004).
\textit{Monte Carlo Methods in Financial Engineering}.
Springer, New York.

\bibitem{HarrisonKreps1979}
Harrison, J.\ M.\ \& Kreps, D.\ M.\ (1979).
\textit{Martingales and arbitrage in multiperiod securities markets}.
Journal of Economic Theory, 20(3), 381--408.

\bibitem{Heston1993}
Heston, S.\ L.\ (1993).
\textit{A closed-form solution for options with stochastic volatility
  with applications to bond and currency options}.
Review of Financial Studies, 6(2), 327--343.

\bibitem{HeEtAl2015}
He, K., Zhang, X., Ren, S., \& Sun, J.\ (2015).
\textit{Delving deep into rectifiers: Surpassing human-level performance
  on ImageNet classification}.
In Proceedings of ICCV, 1026--1034.

\bibitem{Hornik1991}
Hornik, K.\ (1991).
\textit{Approximation capabilities of multilayer feedforward networks}.
Neural Networks, 4(2), 251--257.

\bibitem{KemnaVorst1990}
Kemna, A.\ G.\ Z.\ \& Vorst, A.\ C.\ F.\ (1990).
\textit{A pricing method for options based on average asset values}.
Journal of Banking and Finance, 14(1), 113--129.

\bibitem{KingmaBa2015}
Kingma, D.\ P.\ \& Ba, J.\ (2015).
\textit{Adam: A method for stochastic optimization}.
In International Conference on Learning Representations (ICLR).

\bibitem{LongstaffSchwartz2001}
Longstaff, F.\ A.\ \& Schwartz, E.\ S.\ (2001).
\textit{Valuing American options by simulation: A simple least-squares
  approach}.
Review of Financial Studies, 14(1), 113--147.

\bibitem{LundbergLee2017}
Lundberg, S.\ M.\ \& Lee, S.-I.\ (2017).
\textit{A unified approach to interpreting model predictions}.
Advances in Neural Information Processing Systems (NeurIPS), 30.

\bibitem{Merton1973}
Merton, R.\ C.\ (1973).
\textit{Theory of rational option pricing}.
Bell Journal of Economics and Management Science, 4(1), 141--183.

\bibitem{RumelhartEtAl1986}
Rumelhart, D.\ E., Hinton, G.\ E., \& Williams, R.\ J.\ (1986).
\textit{Learning representations by back-propagating errors}.
Nature, 323(6088), 533--536.

\bibitem{ScornetEtAl2015}
Scornet, E., Biau, G., \& Vert, J.-P.\ (2015).
\textit{Consistency of random forests}.
Annals of Statistics, 43(4), 1716--1741.

\bibitem{SrivastavaEtAl2014}
Srivastava, N., Hinton, G., Krizhevsky, A., Sutskever, I., \&
Salakhutdinov, R.\ (2014).
\textit{Dropout: A simple way to prevent neural networks from
  overfitting}.
Journal of Machine Learning Research, 15(56), 1929--1958.

\end{thebibliography}

\end{document}